% >>> PROPAGATION (Fable 5, Aug 29 2026): at the cosmic-weighting definition, one companion result recorded --- no state-time
% privileged in what the weighting refers to (the no-'next-state' constraint). The Aetherian Nexus is deliberately KEPT OUT of
% this extraction on the author's ruling of this date; its home is the Nexic long paper (N_H, the weighted superposition of all
% mediocre observer-Aethae, equivalently one's earliest possible parent Aethus in Aethic time). See the ledger, R110.
% ============================================================================
% DRAFTING HEADER (strip before deposit)
% Drafted by Claude Fable 5, July 2026, as the fifth extraction of the corpus:
% the anthropic wing's own short paper, drawn from the enhanced Nexic master
% (enhancement passes 1-16: ladder, three explicit proofs, differentia,
% vocabulary apparatus). Author-side card: run the VERIFY sweep against the
% actual volumes per the % comments in the bibliography block below; strip
% those working comments and this header before deposit.
% ============================================================================
\begin{filecontents}[overwrite]{nexic_short_refs.bib}
@misc{Nexus, howpublished={Manuscript}, author={Benander, Ajax}, title={Nexic Reasoning: Defining a Generalized Calculus Over Anthropic Parameters}, year={2024}, month = {12}, url={https://philpapers.org/rec/BENNRD}, doi={10.31219/osf.io/jvkcp}, note={Manuscript; the anthropic wing's master, successive deposits December 2024--2026}}
% VERIFY at merge: sync the AethusBrief entry with its canonical form in the
@misc{AethusSemantics, author={Benander, Ajax}, title={Counterfactual Evaluation by Perturbation of Information States: Deleting the World-Similarity Ordering, Not Refining It}, year={2026}, howpublished={Preprint}, url={https://philpapers.org/rec/BENCEB-2}}
@misc{AethusBrief, author={Benander, Ajax}, title={Aethic Reasoning in Brief: From a Counterfactual Semantics to the Quantum Measurement Problem}, year={2026}, howpublished={Preprint}, url={https://philpapers.org/rec/BENARI-3}, eprint={https://philsci-archive.pitt.edu/id/eprint/30439}, eprinttype={Alternative URL}}
% VERIFY at merge (fields from memory; key matches the corpus canon): Carter,
% in Confrontation of Cosmological Theories with Observational Data, ed. Longair,
% Reidel 1974.
@incollection{carter1974large,
  author = {Carter, Brandon},
  title = {Large Number Coincidences and the Anthropic Principle in Cosmology},
  booktitle = {Confrontation of Cosmological Theories with Observational Data},
  editor = {Longair, M. S.},
  publisher = {Reidel},
  year = {1974}
}
% Referee confirms Routledge 1996 as stated; eyeball once at merge.
@book{leslie1996end,
  author = {Leslie, John},
  title = {The End of the World: The Science and Ethics of Human Extinction},
  publisher = {Routledge},
  year = {1996}
}
% Referee recollection matches 363, 315--319; confirm once at merge.
@article{gott1993implications,
  author = {Gott, J. Richard, III},
  title = {Implications of the Copernican Principle for Our Future Prospects},
  journal = {Nature},
  volume = {363},
  pages = {315--319},
  year = {1993}
}
% VERIFY at merge (publisher details from memory). Note for the transfer sweep:
% the corpus long cites the Rare Earth hypothesis through the key
% 2004IJAsB...3...70E, itself under a lodged VERIFY; this standalone paper
% anchors the hypothesis to the unambiguous book instead --- harmonize the two
% at merge, whichever way the long's VERIFY resolves.
@book{ward2000rare,
  author = {Ward, Peter D. and Brownlee, Donald},
  title = {Rare Earth: Why Complex Life Is Uncommon in the Universe},
  publisher = {Copernicus},
  year = {2000}
}
% VERIFY preferred citation form at merge (working paper).
@misc{hanson1998greatfilter,
  author = {Hanson, Robin},
  title = {The Great Filter --- Are We Almost Past It?},
  year = {1998}
}
% VERIFY at merge (fields from memory): Carter's Royal Society discussion paper.
@article{carter1983anthropic,
  author = {Carter, Brandon},
  title = {The Anthropic Principle and its Implications for Biological Evolution},
  journal = {Philosophical Transactions of the Royal Society of London. Series A},
  volume = {310}, pages = {347--363}, year = {1983}
}
% VERIFY at merge (fields from memory; ApJ 2020 --- confirm volume/article number).
@article{westby2020astrobiological,
  author = {Westby, Tom and Conselice, Christopher J.},
  title = {The Astrobiological Copernican Weak and Strong Limits for Intelligent Life},
  journal = {The Astrophysical Journal}, year = {2020}
}
% R845 (Opus 5.5, Oct 22 2026): three entries checked against arXiv (math/0608592; 0704.2630, Phys. Rev. D 75, 123523;
% 0906.0042, Phys. Rev. D 81, 123524). R845b: Neal under the corpus key Neal2006.
@article{Neal2006,
  author = {Neal, Radford M.},
  title = {Puzzles of anthropic reasoning resolved using full non-indexical conditioning},
  journal = {arXiv preprint}, year = {2006}, note = {arXiv:math/0608592}
}
@article{hartle2007typical,
  author = {Hartle, James B. and Srednicki, Mark},
  title = {Are We Typical?},
  journal = {Physical Review D}, volume = {75}, pages = {123523}, year = {2007}
}
@article{srednicki2010large,
  author = {Srednicki, Mark and Hartle, James},
  title = {Science in a Very Large Universe},
  journal = {Physical Review D}, volume = {81}, pages = {123524}, year = {2010}
}
% VERIFY at merge (fields from memory): arXiv:1806.02404.
@misc{sandberg2018dissolving,
  author = {Sandberg, Anders and Drexler, Eric and Ord, Toby},
  title = {Dissolving the Fermi Paradox},
  note = {arXiv:1806.02404}, year = {2018}
}
@misc{Aethus, howpublished={Manuscript}, author={Benander, Ajax}, title={Aethic Reasoning: A Comprehensive Solution to the Quantum Measurement Problem}, year={2024}, month = {11}, url={https://philpapers.org/rec/BENARA-5}, eprint={https://philsci-archive.pitt.edu/id/eprint/25713}, eprinttype={Alternative URL}, note={Manuscript; the complete development of the framework, successive deposits November 2024--2026}}
@book{bostrom2002anthropic, author = {Bostrom, Nick}, title = {Anthropic Bias: Observation Selection Effects in Science and Philosophy}, publisher = {Routledge}, address = {New York}, year = {2002}}
@article{elga2004defeating,
  author = {Elga, Adam},
  title = {Defeating {Dr.} {Evil} with self-locating belief},
  journal = {Philosophy and Phenomenological Research},
  volume = {69}, number = {2}, pages = {383--396}, year = {2004}
}
@article{snyderbeattie2021timing,
  author = {Snyder-Beattie, Andrew E. and Sandberg, Anders and Drexler, K. Eric and Bonsall, Michael B.},
  title = {The timing of evolutionary transitions suggests intelligent life is rare},
  journal = {Astrobiology},
  volume = {21}, number = {3}, pages = {265--278}, year = {2021}
}
@article{Kaiho2017,
  author  = {Kaiho, Kunio and Oshima, Naga},
  title   = {Site of asteroid impact changed the history of life on {Earth}: the low probability of mass extinction},
  journal = {Scientific Reports},
  volume  = {7}, pages = {14855}, year = {2017},
  doi     = {10.1038/s41598-017-14199-x}
}
@article{mills2025hardsteps,
  author  = {Mills, Daniel B. and Macalady, Jennifer L. and Frank, Adam and Wright, Jason T.},
  title   = {A reassessment of the ``hard-steps'' model for the evolution of intelligent life},
  journal = {Science Advances},
  volume  = {11}, number = {7}, pages = {eads5698}, year = {2025},
  doi     = {10.1126/sciadv.ads5698}
}

@misc{GoldenScience2022,
  author       = {Benander, Ajax},
  title        = {Golden Science},
  year         = {2022},
  month        = {November},
  howpublished = {Self-hosted manuscript},
  url          = {https://aethus.org/aethic-code/files/GoldenScience_Nov2022.pdf},
  note         = {Title page August 2022, the date of the first version of the source; this version created November 2022. Not deposited to a
                  preprint server; cited as a dated developmental record rather than
                  as a source of results.}
}

@article{nesvorny2021dark,
  author  = {Nesvorn{\'y}, David and Bottke, William F. and Marchi, Simone},
  title   = {Dark primitive asteroids account for a large share of {K/Pg}-scale impacts on the {Earth}},
  journal = {Icarus},
  volume  = {368},
  pages   = {114621},
  year    = {2021},
  doi     = {10.1016/j.icarus.2021.114621}
}
\end{filecontents}
\documentclass[11pt]{article}
\usepackage[margin=1.05in]{geometry}
\usepackage{amsmath, amssymb, amsthm}
\usepackage[hidelinks]{hyperref}
\usepackage{microtype}
\usepackage[style=numeric, sorting=none]{biblatex}
\addbibresource{refs.bib} % shared corpus refs, loaded first so they win collisions
\addbibresource{nexic_short_refs.bib}

\newtheorem{theorem}{Theorem}
\newtheorem{principle}{Principle}
\newtheorem{corollary}{Corollary}
\newtheorem{lemma}{Lemma}
\theoremstyle{definition}
\newtheorem{definition}{Definition}

\title{The Copernican Method Against the Copernican Principle:\\ Anthropic Typicality over Partial Information States}
\author{Ajax Benander\thanks{Email: \href{mailto:akb9408@rit.edu}{akb9408@rit.edu}, Website: \url{https://aethus.org}. This paper extracts, as a self-contained argument in anthropic reasoning, the central typicality principle and headline theorem of a larger framework; the companion development \cite{Nexus} carries the natural-historical program built on the principle, which is cited below, not presupposed. The physics-side companions \cite{AethusBrief, Aethus} are likewise cited for the record; their quantum-mechanical claims play no role in the argument.}}
\date{First drafted July 2026\\ {\small This version: September 2026 --- working draft}}
\usepackage{booktabs}
\usepackage{enumitem}

\begin{document}
\maketitle

\begin{abstract}
How can a predeclared record constrain observer abundance under an observer-conditioned sampling model? The Accordance Principle, which governs typicality over partial, observer-indexed information states rather than complete histories and follows from two mediocrity commitments by one application of Bayes' theorem, motivates an observer-conditioned sampling model, under which ordinary probability establishes a finite test. Given a predeclared record of $m$ rare-side realizations at prior below one half, suitable probability bounds, and a stated statistical antecedent imported as a premise, the test rejects an abundance hypothesis $P(H) \ge \varepsilon$ at a level that falls with $m$ and can approach zero. The result is a statistical rejection rather than a contradiction: typicality, assessment under the cosmic weighting over possible histories, and the abundance hypothesis can all hold while the test rejects, with probability at most its level. Applied to the astrobiological Copernican principle, the test addresses an explicit quantitative hypothesis under auxiliary assumptions, and its Rare Earth conclusion holds under cosmic submultiplicativity over a predeclared family of attributes; whether Earth's record meets those conditions is an empirical and selection-modeling question, and the paper's status table does not yet certify its inputs. On the diagnosis proposed, the received reference-class problem is an artifact of reasoning over completed histories, with the choice of sampling measure and granularity remaining in its place.
\end{abstract}

\section{Introduction}
\label{sec:intro}

% R845 (Opus 5.5, Oct 22 2026), with the referee's report of that date: (P1) common ground for most parties, not all
% (Hartle and Srednicki); it read "for every party".
Anthropic reasoning trades on typicality: whatever else is uncertain, one should not expect to be a wild exception among observers. That commitment, in some form, is common ground for most parties to the standard debates, though not for all \cite{hartle2007typical}. This paper argues that the commitment, taken together with two further theses each of which its practitioners standardly also hold, generates a statistical test against their conjunction whose level falls without bound in the number of predeclared rare-side realizations, conditional on one natural-historical result imported from a companion development and flagged where it enters.

\textit{The triad.} The three members of the triad are: \emph{self-mediocrity}, the thesis that one's own situation is typical among observer-situations under the operative weighting; \emph{cosmic-mediocrity}, the thesis that the operative weighting is the cosmic one---occurrence-weighted over the space of possible histories---rather than a gerrymandered restriction of it; and the \emph{astrobiological Copernican principle}, the thesis that observer-bearing environments like Earth's are typical among planetary environments, so that complex life should be common --- a thesis with genuine currency in the modern estimate literature \cite{westby2020astrobiological}, not a construction of this paper's.



\textit{The engine.} The engine of the argument is a principle governing typicality, stated and derived in Section~\ref{sec:principle}: the Accordance Principle, on which the prior improbability of a realized natural-historical attribute must be compensated---accorded---by that attribute's relevance to the production of observers. The principle is not a new axiom grafted onto anthropic reasoning; it is one application of Bayes' theorem away from the two mediocrity commitments, once the object those commitments range over is corrected.

\textit{The plan.} The correction is the subject of Section~\ref{sec:relocations}: the framework developed here differs from observation-selection theory as standardly practiced by two relocations, and the relocations carry everything. Section~\ref{sec:theorem} then proves that a finite, predeclared family of attributes supplies a level-$\delta$ test against the triad's conjunction, with the one empirical link in the chain---from the principle to the Rare Earth hypothesis---derived here, its structural half as a theorem and its empirical antecedent argued as the weakest input the theorem admits. Section~\ref{sec:exits} prices the three exits and argues for surrendering the third member. Section~\ref{sec:objections} answers the four objections nearest to hand, and Section~\ref{sec:fence} states what this paper does not claim.

\noindent\emph{What the paper establishes, and what it opens.} The paper's results are of two kinds, held to different standards. Its theorems---the rarity bound, the realized-cell bound, the finite Rare Earth corollary, and the statistical trilemma---are mathematics: each holds for any inputs satisfying its hypotheses, and none needs data to be true. Its statements about our own history are outputs of that calculus under stated inputs: the predeclared family of attributes, the cosmic priors of the sides a record realizes, and the product condition on those priors, (A1) below. Where the paper states a figure for Earth, it is the calculus run on inputs it names, offered as an illustration and not as a measurement. Estimating those inputs is not a debt the framework owes before it counts; it is the program the framework opens. As empirical data accumulate and the couplings among natural-historical attributes are worked out, natural history stops being a tally of candidate improbable events and becomes a topography: productivities falling away from the realized history along measurable gradients, contrast spreads that can be estimated, and structure whose presence the recorded falsifiers decide. On that map, Nexic inference becomes optimization over a high-dimensional manifold of histories, and the calculus given here is the part of the map that holds at every resolution, as field equations hold whatever the metric. The long development maps that program \cite{Nexus}.

\section{The argument in ordinary Bayesian terms}
\label{sec:plain}

The apparatus this paper runs on is unfamiliar, and the argument it licenses is not. Before the framework is introduced it is worth seeing the whole thing once in plain conditional-probability terms, so that a reader can tell which part of what follows is machinery and which part is the point.

Start from the recognizable Copernican statement: \emph{among all the situations in which an observer like us could find themselves, ours should not be a special one}. Take an attribute $A$ of our natural history---that a giant impact struck hydrocarbon-rich sediment, say---with alternative $A'$, and let $H$ be the production of an observer. Bayes' theorem relates the observer-conditioned odds to the prior odds:
\begin{equation}
\label{eq:plain-bayes}
\frac{P(A' \mid H)}{P(A \mid H)} \;=\; \frac{P(H \mid A')}{P(H \mid A)} \cdot \frac{P(A')}{P(A)} .
\end{equation}
% R836 (Opus 5.5, Oct 21 2026), at the author's direction: P(H|A')/P(H|A) named "the canonical ratio" here and used by name below.
Call the middle factor, $P(H \mid A')/P(H \mid A)$, the alternative's observer-productivity relative to the realized side's, the \emph{canonical ratio}; it is the quantity everything below bounds. Now read typicality as a constraint on the left-hand side: if our situation is not special among observer-situations, the realized side of the partition should not carry a small share of the observer-conditioned weight. Demanding that share be at least $\tfrac12$ makes the left-hand side at most one, and \eqref{eq:plain-bayes} becomes
\[ \frac{P(H \mid A')}{P(H \mid A)} \;\lesssim\; \frac{P(A)}{P(A')} , \]
which is the Accordance Principle, the canonical ratio at most the prior odds of the realized side, with $\lesssim$ marking that typicality makes the demand more often than not, so that the bound is a median and not an absolute (\autoref{sec:principle}): \emph{a realized attribute's prior improbability must be compensated by its relevance to producing observers}. If the impact site was a one-in-eight shot and yet was realized, then either the alternative sites were correspondingly worse at producing observers, or our situation is atypical. Nothing in that derivation is exotic. It is one application of Bayes to one attribute, with typicality supplying the inequality. The application has two stages. The first is the inequality itself, which registers a realized attribute of low conditioned mass as a fluke that must be compensated. The second is \emph{teleological discharge}: identifying, from outside context or further inference, where the fluke helps, the attribute's contribution to observer-production that the first stage only says must exist; the two can be separated by a long interval, the fluke standing undischarged until its context arrives. ``Teleological'' is the Nexic sense only, an inferential aid exactly as Darwinian teleology is, survivorship bias and not purpose, the distinction being this paper's standing disclaimer. In the impact example the site's improbability is the first stage and the cleared niche for the mammalian radiation its discharge.
% R871 (Opus 5.5, Oct 27 2026), at the author's direction (Oct 27: \lesssim for the effective forms, the heuristic median, as the
% running convention): the display, which read "\le", written with \lesssim, and the clause on it added.

Two things the rest of the paper adds to it. The first is a \emph{tail}: ``not special'' is a demand one can make at any strength, and demanding a share of at least $1/(1+\kappa)$ rather than $\tfrac12$ turns the display into a family of bounds indexed by $\kappa$, each with its own confidence, which is what makes the principle quantitative rather than a slogan. The second is the object typicality is asserted \emph{of}, and this is where the paper departs from the tradition.

One thing about the derivation just run is worth noticing before it is generalized: Bayes' theorem does not care which of $A$ and $H$ came first. Inverting a conditional privileges neither event as the earlier, and the framework's tense-free reading of openness is the property Bayes already respected without being asked to---as Maxwell's equations respected special relativity before there was one to respect. The framework supplies the object on which that indifference becomes usable; it does not supply the indifference. The recognizable statement above quantifies over observers---``among all the situations in which an observer could find themselves.'' That phrasing is the source of the tradition's central difficulty: which observers, in which reference class, counted how? The question has no stable answer, and the literature's long dispute between self-sampling and self-indication is a dispute about it, (\autoref{sec:relocations}). This paper does not answer the question. It asserts typicality of something else: not of the observer's position among observers, but of \emph{the weight an attribute's realized value carries within its own partition}. The draw is over the attribute's answer space, the class is the attribute's own alternatives, and the record supplies both. No observer is counted anywhere in the derivation, which is why the reference-class question is relocated to the measure, not answered, and why the self-sampling dispute is one the paper is not party to.

That is the whole of the exchange the framework asks a reader to make. In return for reasoning over partial information states---attributes of a history with the rest left open---rather than over completed histories with observers distributed through them, one gets a typicality principle that applies attribute by attribute, needs no reference class, and yields a bound that can be computed from published priors. Everything after this section is that exchange made precise: what a partial state is, what it means for a weight to be typical, and what follows when the principle is applied to many attributes at once.

\section{The framework in brief}
\label{sec:framework}

The object of reasoning throughout is a \emph{partial information state}:\footnote{A naming note, recorded once: the fuller developments give the partial states a proper name, \emph{Aethus}, plural \emph{Aethae} \cite{Nexus, Aethus}.} a partial specification of reality, individuated by what it states and closed under what its content entails, with everything unstated left genuinely blank rather than secretly filled in. An observer's epistemic situation is such a state; so is any hypothetical stage of a natural history, specified up to some set of attributes and no further. The discipline the object imposes is the discipline of the blank: an unstated attribute is not a hidden variable awaiting revelation but a genuine absence, and no inference below is permitted to fill a blank as a step---every constraint must bind on the stated content alone. Two definitions fix the usage.

\begin{definition}[Aethus, or partial information state]
An \emph{Aethus} (plural \emph{Aethae}) is a partial specification of the way things go --- some attributes stated, all others blank --- individuated by its stated content and closed under entailment from that content. No completed history is presupposed as the object the state partially describes; the Aethus itself is the unit of specification, and the name marks that: no completed history sits silently underneath the object, which is what distinguishes it from an incomplete Bayesian description of a history that is there. ``State'' abbreviates ``Aethus'' below.
\end{definition}
\noindent The calculus this paper develops over Aethae---anthropic typicality, the Accordance principle, and the bounds that follow---is \emph{Nexic reasoning}, one branch of \emph{Aethic reasoning}, the general method whose objects are Aethae and whose operations preserve their partiality rather than completing them: Aethic reasoning contains Nexic reasoning as the typicality calculus over its objects.

\begin{definition}[Weighting; the cosmic weighting]
A weighting is a probability measure over partial information states. The cosmic weighting is the weighting induced by occurrence: states weighted by the measure of ways the world can realize them, with no restriction to any privileged subfamily and no conditioning silently absorbed into the measure. The measure is defined on the states directly; it need not arise as the marginal or pushforward of any measure on completed histories. This is the weighting the companion names the \emph{Cosmic Nexus} \cite{Nexus}. One thing about it is fixed by the fuller developments and recorded here since it bears on how the measure is read: \emph{no state-time is privileged in what the weighting refers to}: the weight of a far descendant of the present state is not assembled by chaining through intermediate states---a far descendant is a direct weighted refinement of the present state, its weight fixed now---so the weighting is as informative about the state one may hold ten years hence as about the one held a minute hence, and ``the next state'' is not a thing the measure singles out \cite{Aethus}.
\end{definition}

Nothing in this paper requires the full ontology of the companion frameworks---in particular, nothing here turns on the physics-side thesis that completed histories are ill-formed rather than merely dispensable \cite{Aethus}---but one consequence of that ontology is used as a working posture: no step below ever requires completing a partial state into a total history, and every constraint is stated attribute by attribute.

Two mediocrity commitments are now statable cleanly. \emph{Self-mediocrity} (P1): one's own situation is typical under the operative weighting restricted to observer-situations---formally, each interrogated attribute's realized value is distributed as a draw from that observer-conditioned weighting over its answer space, the Copernican draw stated formally in Section~\ref{sec:principle}---in the companion's intuitive gloss, ``if it is rare to everybody, it should almost always be rare for you as well'' \cite{Nexus}. \emph{Cosmic-mediocrity} (P2): the operative weighting is the cosmic one. The second commitment is what forbids gerrymandering: typicality claims are cheap if the reference weighting may be tailored after the fact, and the cosmic weighting is the non-tailored choice. Two things are distinct, since running them together would overstate what P2 delivers. That a weighting is \emph{not gerrymandered} is not the same as its being \emph{uniquely determined}: P2 rules out weightings selected to suit the attribute under interrogation, and it does not single out one measure from among the admissible ones. What the framework claims is the weaker robustness that the results actually use---the Accordance display is a statement about a \emph{ratio}, so any reweighting that does not differentially target the $A/A'$ pair leaves it unchanged, and two admissible cosmic weightings disagreeing only in ways orthogonal to that pair yield the same bound. Reweightings that do differentially target the pair are what P2 forbids. One further conditionality belongs here, since a standalone reader will ask for it. The paper does not construct the measure. It assumes the existence of a coherent probability measure over partial states satisfying the stated properties---in particular, that the displayed $P(A)$, $P(H \mid A)$ and their ratios inherit the ordinary probability identities---and every result is conditional on that existence. Supplying the measurable structure on the overlapping family of partial states is the framework's business rather than this paper's, and a reader who doubts that such a measure exists doubts the framework at its root rather than any step taken here. This is a narrower guarantee than uniqueness, and the paper does not claim more: choosing among admissible cosmic measures remains the ordinary Bayesian burden of choosing a prior, which the framework inherits rather than dissolves, and which it holds to be a different and lesser difficulty than choosing an observer reference class, (\autoref{sec:relocations}).

\section{The Accordance Principle}
\label{sec:principle}

\begin{principle}[The Accordance Principle]\footnotemark
\label{prin:accordance}
For a natural-historical attribute $A$ with alternative $A'$ --- canonically its complement, so that $P(A') = 1 - P(A)$ and the boxed display matches the companion's \cite{Nexus} --- and $H$ the production of an observer, every probability relative to a declared background state in which the interrogated attribute is blank, the draw taken over that attribute's answer space,
\begin{equation}
\label{eq:exceedance}
\boxed{\;\Pr\!\left[\, \frac{P(H \mid A')}{P(H \mid A)} \;\ge\; \kappa\, \frac{P(A)}{P(A')} \,\right] \;\le\; \frac{1}{1+\kappa}, \qquad \kappa \ge 1,\;}
\end{equation}
the probability taken over the realization, and the event strict at $\kappa = 1$; its effective form is its median, at $\kappa = 1$,
\begin{equation}
\label{eq:accordance}
\frac{P(H \mid A')}{P(H \mid A)} \;\lesssim\; \frac{P(A)}{P(A')},
\end{equation}
read attribute by attribute over a partial state, with the canonical ratio on the left: the prior improbability of a realized attribute must be accorded---compensated---by that attribute's relevance to the production of observers.
\end{principle}
% R873 (Opus 5.5, Oct 27 2026), at the author's direction (Oct 27: the right things boxed, the general form before the
% special overall): the principle stated with its general form first, the boxed exceedance display moved here from the
% end of the derivation's third step, and its effective form, the median, second; "the display matches the companion's
% boxed form" read "the boxed display matches the companion's", the companion now boxing the exceedance form.
\footnotetext{A dating and origin note, retained as provenance, with the principle's retrospective reinterpretation. The original motivation (2019--2020, predating the framework): imagine an infinite universe in which all possibilities happen somewhere, so that the number of ``Earths'' on which a given event occurs is proportional to that event's probability of occurring; the Accordance Principle then deduces as simply as dropping oneself on a random such Earth. The retrospectively remarkable feature is that this is accidentally an entirely Aethic argument --- Aethic precisely because the discriminating attribute is valued as \emph{non-exhaustive}, everything neither specified by nor deducible from the stated attribute standing in Aethic superposition relative to it --- unrecognized at the time under a still-classical paradigm. In the Aethic writeup \cite{AethusSemantics}, probability is valued over superposition objects --- Aethae --- rather than over determinate classical worlds, so the standard anthropic dictum ``the probability of being born somewhere is proportional to the population size'' becomes ``the probability of existing somewhere is proportional to the population size times the Aethic weight.'' Qualitatively, a population turns \emph{translucent} under partial weight: within the linear combination it matters how partial a population's weight is, the random drawing occurring at the interface between Aethic realities. The principle as used here then assumes each eventuality carries the same population size --- eventualities being Aethae, they hold multiple determinate population sizes in disagreeing superposition under the extra unknown information, hence inferring in the ``do observers exist'' case of the Accordance Principle that these differences might average out to common similarity. If, however, they do not, then simply by switching ``do observers exist'' to an Aethus holding more attributes targeted at a more individual person such as one's own Aethus, one attains that a smaller population becomes interpretable simply as a lower probability of generating such a person in particular. The derivation's own history completes the note, dictated for the record in August 2026, with one calibration the author insists on: the construction was seeded by popular exposure to the Doomsday argument --- heard of online before 2019, no paper of the tradition read --- and it diverged from the literature precisely because the reaction to that much was phenomenological rather than textual. The original construction was an infinite-cosmos \emph{Earth count} --- the number of Earths satisfying $A$ against a counterfactual $A'$, relative to the branch point $B = A \cup A'$, their ratio read directly as the probabilities' ratio --- from which the Copernican draw, and the principle, followed as consequences. By 2021 the Earth count and timeline-counterfactuals stood intuitively unified; a 2022 intermediate replaced the census with a \emph{mass of timelines} under each probability, transitional between the discrete count and value over states; and the 2022--2023 step re-valued the premises over indexed partial states --- past light-cone attributes able to roll from scratch, the born-once-on-a-random-Earth picture retired --- which is what turned a striking property of an imagined cosmos into a deduction from the relativity of determinacy.}

\begin{proof}[Derivation]
Self-mediocrity, applied to the realized attribute, says the realized value should be typical under observer-conditioning: $P(A' \mid H) \lesssim P(A \mid H)$, since if the alternative were both likelier a priori and at least as productive of observers, finding the realized value would make one's situation atypical among observer-situations. Cosmic-mediocrity fixes the weighting in which these probabilities are taken. Bayes' theorem gives
\[
\frac{P(A' \mid H)}{P(A \mid H)} \;=\; \frac{P(H \mid A')}{P(H \mid A)} \cdot \frac{P(A')}{P(A)},
\]
and the typicality condition is therefore the effective form \eqref{eq:accordance}; the general form is derived below. The content is direct: where the realized attribute was the a priori unlikely member of the pair, the constraint forces the canonical ratio below one, the alternative having been the less observer-productive member---improbability on the path to an observer must be bought by observer-relevance.
\end{proof}

The relation $\lesssim$ carries the principle's modality, and it can be given a definite meaning rather than left as a tolerance. The formal content of self-mediocrity is the \emph{Copernican draw}, and on that formalization the bound needs no tolerance parameter: it is the mediocrity principle read at a threshold and rewritten by Bayes' theorem. The derivation runs in three steps.

\textit{Step 1: the mediocrity principle.} Drop a point at random on the unit interval and let $a$ be where it landed. The principle is one inequality on the span $[0,a]$:
\begin{equation}
\label{eq:mediocrity}
\Pr\!\big[\, m[0,a]\,/\,m[0,1] \;\le\; w \,\big] \;\le\; w, \qquad w \in [0,1].
\end{equation}
For the span one landed in to take up no more than a proportion $w$ of the whole happens no more than a proportion $w$ of the time. The inequality is the principle and the equality is its idealization: on an unconstrained interval the landing quantile is uniform and the bound is attained---the probability integral transform---but an Aethic partition is not an unconstrained interval, since the source framework's pruning removes cells, and where the cells removed are light ones the event that the realized span is that light loses measure and the probability falls strictly below $w$.\footnote{The $\Pr$ here, and every $\Pr$ descending from it below, is the source framework's \emph{Aethic frequentist probability} \cite{Aethus}, a proportion over the declared class of draws, and not an Aethic weight; the outer operator of the Accordance bound is frequentist for the same reason, and the derivation never changes what $\Pr$ is, only what it ranges over.}

\textit{Step 2: the Aethic Copernican Draw.} The draw is \eqref{eq:mediocrity} applied to an Aethic partition. Let $H$ be the Aethus stating only that observers are produced, let $H^{\ast}$ be the reasoner's own Aethus, and interpose two Aethae between them,
\begin{equation}
\label{eq:chain}
H^{\ast} \;\sqsubseteq\; H^{-} \;\sqsubseteq\; H^{+} \;\sqsubseteq\; H,
\end{equation}
each a refinement of the next in the intersection sense. The draw asserts that $H^{-}$ is mediocre in $H^{+}$---the two middle Aethae, not $H^{\ast}$ in $H$---so that the middle may be chosen anywhere the chain's validity permits. The chain's ends do different work, and inclusion does not transport the sampling law: with $P(A) = 10^{-3}$, $H = A$ and $H^{+} = \Omega$, an observer sampled conditional on $H$ sees $A$ with certainty while $P(H^{+}) = 1$, so applying a rarity bound to an ancestor of $H$ needs a screening or transport condition beyond the chain relation itself. At the top, $H$ supplies the observer-conditioning: every Aethus in \eqref{eq:chain} lies inside it, so every probability below is observer-conditioned by construction and no separate tilt is needed. At the bottom, $H^{\ast}$ constrains the tower: nothing above it may state what the reasoner's own records contradict. The attribute and its alternative are defined by the two spans,
\begin{equation}
\label{eq:attribute-def}
H^{-} \;\equiv_{\mathcal{V}}\; A\,H^{+}, \qquad H^{+} - H^{-} \;\equiv_{\mathcal{V}}\; A'\,H^{+},
\end{equation}
where $\equiv_{\mathcal{V}}$ is validity equivalence rather than the stricter identity of stated content, and which in the set notation of ordinary intuition read $H^{-} = A \cap H^{+}$ and $H^{+} \setminus H^{-} = A' \cap H^{+}$. The spans are Aethae rather than intervals of time, so by tense-free blankness $A$ may be past, present, or future. Taking $H^{+} = H$ and substituting $m[0,a] \mapsto P(H \cap A)$, $m[0,1] \mapsto P(H)$, and $w \mapsto 1/(1+\kappa)$ into \eqref{eq:mediocrity} gives
\begin{equation}
\label{eq:draw}
\Pr\!\big[\, P(A \mid H) \;\le\; \tfrac{1}{1+\kappa} \,\big] \;\le\; \frac{1}{1+\kappa}, \qquad \kappa \ge 1,
\end{equation}
with the event strict at $\kappa = 1$; the range of $\kappa$ is where the substitution earns its keep, so it is derived rather than assumed. The continuous span $[0,a]$ of \eqref{eq:mediocrity} has a uniform proportion; the Aethic partition is not a continuum but two cells, $A$ and $A'$, of fixed observer-conditioned weights $q$ and $1-q$, and the realized cell's weight $Q = P(A \mid H)$ takes the value $q$ with probability $q$ and $1-q$ with probability $1-q$. Then $\Pr[Q \le w]$ is the total weight of the cells of weight at most $w$, which is at most $w$ exactly when at most one cell can qualify---that is, when $w < \tfrac12$, since two cells of weight at most $w \ge \tfrac12$ can together carry all the mass. At $w = \tfrac12$ the inequality holds in the strict form $\Pr[Q < \tfrac12] \le \tfrac12$, failing only at an exact tie; for $w > \tfrac12$ it fails outright wherever both cells are light enough, as at $q = \tfrac12$, where $Q = \tfrac12$ with certainty and $\Pr[Q \le 0.6] = 1$. So the draw's super-uniformity is a statement about the lower tail, $\kappa \ge 1$, with $\kappa = 1$ in strict form, and nothing below uses more.\footnote{An earlier statement of this derivation claimed \eqref{eq:draw} for every $\kappa > 0$, transferring the continuous interval's full uniformity to the two-cell partition, and the comprehensive development's own record shows the $\kappa \ge 1$ constraint being dropped on that reasoning \cite{Aethus}. The two-cell counterexample decides the question; the constraint stands with its reason, and the many-attribute bound is derived from the lower-tail statement in \S\ref{sec:rarity-bound}.}

\textit{Step 3: Bayes' theorem.} The event in \eqref{eq:draw} is rewritten in two lines. Since $P(A \mid H) + P(A' \mid H) = 1$,
\begin{equation}
\label{eq:complement}
P(A \mid H) \;\le\; \tfrac{1}{1+\kappa} \quad\Longleftrightarrow\quad P(A' \mid H) \;\ge\; \kappa\, P(A \mid H),
\end{equation}
and by Bayes' theorem, $P(X \mid H) = P(H \mid X)\,P(X)/P(H)$,
\begin{equation}
\label{eq:bayes}
P(A' \mid H) \;\ge\; \kappa\, P(A \mid H) \quad\Longleftrightarrow\quad \frac{P(H \mid A')}{P(H \mid A)} \;\ge\; \kappa\, \frac{P(A)}{P(A')}.
\end{equation}
Substituting \eqref{eq:complement}--\eqref{eq:bayes} into \eqref{eq:draw} gives the principle's general form, \eqref{eq:exceedance}, the probability taken over the realization, the event's left-hand side being the canonical ratio and its inequality strict at $\kappa = 1$. In the idiom of surprise: to find oneself on a statistically surprising branch is evidence that the branch was less surprising from the standpoint of the process that produces observers than from the standpoint of the prior, and \eqref{eq:exceedance} says by how much.
% R873 (Opus 5.5, Oct 27 2026), at the author's direction: the boxed display that stood here, after "gives the
% principle:", moved into the principle's statement; "the display says by how much" read "\eqref{eq:exceedance} says
% by how much". In the derivation after the statement, "the typicality condition is therefore the displayed bound"
% read "the effective form \eqref{eq:accordance}; the general form is derived below".

\textit{Three readings of the bound.} Write $R = P(A' \mid H)/P(A \mid H)$ for the posterior odds against the realized attribute, the canonical ratio divided by the prior odds $P(A)/P(A')$, and $Q = P(A \mid H)$ for its proportion, so that $R = (1-Q)/Q$. First, the \emph{quantile ladder}: taking $\kappa = (1-p)/p$ in \eqref{eq:exceedance} gives
\begin{equation}
\label{eq:ladder}
\Pr\!\left[\, R \;\ge\; \tfrac{1-p}{p} \,\right] \;\le\; p, \qquad p \in (0, \tfrac12),
\end{equation}
% R845: E[R] = 1 for two-cell draws and any mixture of them; the uniform Q no such mixture but the envelope (it led with
% "Where Q is uniform ... diverges").
so the $(1-p)$-quantile of $R$ is at most $(1-p)/p$: the ninetieth percentile at most $9$, the ninety-ninth at most $99$, and the median at most $1$ in the strict form $\Pr[R > 1] \le \tfrac12$, which is the $p = \tfrac12$ rung read with strict inequality, the only form the two-cell draw supports there. Display~\eqref{eq:accordance}'s $\lesssim$ is this ladder's middle rung, $p = \tfrac12$, and it asserts that the \emph{median} of $R$, the canonical ratio divided by the prior odds, is at most one. The other rungs keep the exceedance form, the default statement of every rung, whose $\Pr[\cdot]$ carries the probability and inside which the inequality is written exactly; the tendency symbol is reserved for what would otherwise be written as a bare inequality, above all the median. Where another rung is abbreviated so, a subscript gives its failure probability, $x \lesssim_q y$ asserting that $x \le y$ fails with probability at most $q$, fixed in advance, so that a bare $\lesssim$ is $\lesssim_{50\%}$, and $x \gtrsim_q y$ is $y \lesssim_q x$. The statistic is the median and not the mean. A nondegenerate fixed binary partition has $\mathbb{E}[R] = 1$ and $\mathbb{E}[1/Q] = 2$ for every $q$, and any mixture of such draws keeps $\mathbb{E}[R] = 1$; a uniform $Q$, under which $\mathbb{E}[R] = \int_0^1 (1-q)/q\,dq$ diverges, is no such mixture but the envelope that attains the ladder. A mean ratio of one says nothing about the realized ratio's typical size, which is what the median rung constrains.
% R872 (Opus 5.5, Oct 27 2026), at the author's direction (Oct 27: the Pr[] form the default throughout, \lesssim used
% only where an inequality would otherwise stand bare, as at the median or effective form): the convention, added in
% R871, restated with the exceedance form first and the subscripted symbol an abbreviation only; R871's statement of
% the ladder in the subscripted notation withdrawn, the ladder standing in its Pr[] form, eq:ladder.

% R836 (Opus 5.5, Oct 21 2026), at the author's direction: the principle's percentile scale for kappa, as a table.
% R843 (Opus 5.5, Oct 22 2026), with the referee's third report on Optimal Earth: kappa_beta a threshold at a
% guaranteed level (was "at the beta-th percentile of the realization"); R's own quantile lies at or below it, as the
% ladder above says. The table's header "Level beta" and "Canonical bound" (they read "Percentile beta", "Canonical ratio
% at most"); "level one half" for "the 50th percentile".
\textit{The percentile scale.} The ladder is a percentile scale for $\kappa$. At guaranteed level $\beta$ the canonical ratio is at most $\kappa_\beta = \beta/(1-\beta)$ times the prior odds, the bound failing with probability at most $1 - \beta = 1/(1+\kappa_\beta)$ (Table~\ref{tab:kappa-scale}). The scale's check is its middle: $\kappa = 1$ sits at level one half. It also begins there. Below the median a tighter bound would need the realized side to outweigh its alternative by a margin, which an attribute whose two sides carry nearly equal observer-weight defeats; below level one half the bound therefore stays at the median's, $\kappa = 1$. Where the prior $P(A)$ is itself known only up to a confidence distribution, the bound is read on two percentiles, this scale's and the prior's, or on one, the canonical quantiles of \eqref{eq:canonical-quantile} below, the complement of the level of \eqref{eq:level-reading} averaged over the prior's confidence distribution with no credit below $\kappa = 1$, a weighted calibration index, not a guarantee at the unknown prior.

\begin{table}[ht]
\centering
\small
\caption{The Accordance Principle's percentile scale. At guaranteed level $\beta$ the canonical ratio $P(H \mid A')/P(H \mid A)$ is at most $\kappa_\beta = \beta/(1-\beta)$ times the prior odds $P(A)/P(A')$, the bound failing with probability at most $1/(1+\kappa_\beta)$; the last column reads it for an attribute of prior $0.01$. The scale begins at level one half, $\kappa = 1$ in strict form; below it the bound stays at the median's. For a given prior the failure probability is the canonical quantile of \eqref{eq:canonical-quantile}.}
\label{tab:kappa-scale}
\begin{tabular}{@{}cccc@{}}
\toprule
Level $\beta$ & $\kappa_\beta$ & Fails with probability at most & Canonical bound, prior $0.01$ \\
\midrule
$50$ & $1$ & $1/2$ & $0.0101$ \\
$60$ & $1.5$ & $2/5$ & $0.0152$ \\
$75$ & $3$ & $1/4$ & $0.0303$ \\
$80$ & $4$ & $1/5$ & $0.0404$ \\
$90$ & $9$ & $1/10$ & $0.0909$ \\
$95$ & $19$ & $1/20$ & $0.192$ \\
$99$ & $99$ & $1/100$ & $1.00$ \\
\bottomrule
\end{tabular}
\end{table}

Second, a substitution the theorem does not license. Taking $\kappa = (1 - P(A))/P(A)$ would give $\Pr[P(H \mid A') \ge P(H \mid A)] \le P(A)$, a bound on the prior alone. The substitution is illegitimate, because $A$ is the realized member and $P(A)$ therefore varies with the realization, making $\kappa$ a function of the draw; a tail bound proved for each fixed $\kappa$ does not license a $\kappa$ chosen from the outcome. The decisive case is independence: there, $P(H \mid A') = P(H \mid A)$ holds with certainty, so the event has probability one and the bound would read $1 \le P(A)$. The null-case corollary below says correctly that independence returns nothing, and the two could not both stand; the prior bound is withdrawn. What a $\kappa$ read off the prior does license is a level. Fix a factor $c > 0$ before the realization, let $\kappa_c = c\,P(A')/P(A)$, the $\kappa$ that puts the bound at $c$, and let $\lambda_c = \kappa_c/(1+\kappa_c)$ when $\kappa_c \ge 1$, with no level otherwise. Then for every $\beta \in [\tfrac12, 1)$
\begin{equation}
\label{eq:level-reading}
\Pr\!\left[\, \frac{P(H \mid A')}{P(H \mid A)} \;>\; c \ \text{ and } \ \lambda_c \;\ge\; \beta \,\right] \;\le\; 1 - \beta,
\end{equation}
since by \eqref{eq:complement}--\eqref{eq:bayes} the canonical ratio exceeds $c$ exactly when $Q < 1/(1+\kappa_c) = 1 - \lambda_c$, so that the event lies inside $\{Q < 1 - \beta\}$, a fixed threshold at or below one half, which \eqref{eq:draw} bounds by $1 - \beta$, in strict form at one half. The level is valid however it depends on the realized prior; at $c = 1$ it is $1 - P(A)$ where $P(A) \le \tfrac12$, with no level otherwise; and the event is strict, so that for $c \ge 1$ it is empty at independence, where the canonical ratio is exactly one.

% R840 (Opus 5.5, Oct 22 2026), at the author's direction: canonical quantiles, as the long paper defines them; R840b,
% reoriented the same day, the canonical quantile of c being Pr[V > c].
% R842: the confidence-weighted level named a weighted calibration index, as the referee's report on Optimal Earth
% asks (it read "not as a guarantee at the unknown prior").
% R846 (Fable 5.1, Oct 23 2026), at the author's direction, the additions' formatting: the two laws displayed (eq:two-laws),
% the closure operator displayed (eq:closure-operator); to hold 31 pages, the fair-coin aside and the placement footnote's
% "now checkable line by line" clause dropped.
% R843: r_q named the canonical bound, a threshold of V (was "the canonical ratio at canonical quantile q"); V read at a
% guaranteed level drawn at random; the realization guarantee, at the realized side's prior odds, and the auxiliary label
% kept apart.
\textit{Canonical quantiles.} The ladder and the level reading are one scale read in two directions. For $q \in (0,1)$ put $\kappa(q) = \max\{1, q/(1-q)\}$, and let $V = \kappa(U)\, P(A)/P(A')$ with $U$ uniform on $(0,1)$, the principle's bound read at a guaranteed level drawn at random, an auxiliary variable. The \emph{canonical quantile} of a value $c$ is $\Pr[V > c]$, which falls as $c$ rises, and the \emph{canonical bound} $r_q$ at canonical quantile $q$, a threshold of $V$ falling as $q$ rises, is the least $c$ with $\Pr[V > c] \le q$: a higher canonical quantile asks more of the realized side. For a given prior,
\begin{equation}
\label{eq:canonical-quantile}
r_q \;=\; \kappa(1-q)\, \frac{P(A)}{P(A')}, \qquad \Pr[\,V > c\,] \;=\; \frac{P(A)/P(A')}{c + P(A)/P(A')} \quad \text{for } c \;\ge\; \frac{P(A)}{P(A')},
\end{equation}
and $1$ below: the canonical quantile of $c$ is one minus the level $\lambda_c$ of \eqref{eq:level-reading}, and $r_q$ is its inverse. Two statements concern two different laws,
\begin{equation}
\label{eq:two-laws}
\Pr_{\mathrm{aux}}[\, V \le r_q \,] \;\ge\; 1-q \qquad\text{and}\qquad \Pr\!\Big[\, \tfrac{P(H \mid A')}{P(H \mid A)} \;>\; \kappa(1-q)\, \tfrac{P(A)}{P(A')} \,\Big] \;\le\; q,
\end{equation}
the first the label, under the law of $U$; the second the ladder \eqref{eq:ladder}, over the realization, for $q \in (0, \tfrac12]$ fixed in advance and with its threshold set by the realized side's prior odds, so that for a given prior the canonical quantile is the failure probability of Table~\ref{tab:kappa-scale}, and above one half $r_q$ is the median's bound. Where $P(A)$ is known only through a confidence distribution, it is drawn from that distribution independently of $U$, and the canonical quantiles are those of the resulting $V$: the canonical quantile of $c$ is then one minus the level of \eqref{eq:level-reading} averaged over the distribution, the label of \eqref{eq:two-laws} is a weighted calibration index, not a posterior or a coverage guarantee at the unknown prior, and above one half $r_q$ leaves the median's bound only through the prior's spread.

Third, the \emph{majority reading}, which is the phenomenology the principle is named for and is a corollary of the inequality's shape rather than a separate posit. By the survival identity and \eqref{eq:mediocrity},
\begin{equation}
\label{eq:majority}
\mathbb{E}[Q] \;=\; q^{2} + (1-q)^{2} \;=\; \tfrac12 + 2\big(q - \tfrac12\big)^{2} \;\ge\; \tfrac12,
\end{equation}
with equality exactly at $q = \tfrac12$, the fair-coin attribute. The proof is the two-cell computation directly---$Q$ is $q$ with probability $q$ and $1-q$ with probability $1-q$---and not the survival-function integral $\int_0^1 (1 - \Pr[Q \le w])\,dw \ge \int_0^1 (1-w)\,dw$, which would need $\Pr[Q \le w] \le w$ on all of $[0,1]$ and has it only below $\tfrac12$. The realized attribute is \emph{expected} to lie in the greater span; its proportion is super-uniform in its lower tail, in the sense that a conservative $p$-value is super-uniform under its null. The expectation is available here and not in the ratio coordinate because $Q$ is bounded in $[0,1]$ where $R$ is not; the bounded transform $Q = 1/(1+R)$ is what carries the statement across.

Fourth, the \emph{null case}. Let $A$ be statistically independent of $H$, so that $P(H \mid A) = P(H \mid A')$. Then the canonical ratio in \eqref{eq:exceedance} is identically $1$ and the event becomes, for every $\kappa > 1$,
\begin{equation}
\label{eq:null-case}
\Big\{\, 1 \;\ge\; \kappa\, \frac{P(A)}{P(A')} \,\Big\} \;=\; \Big\{\, Q \;\le\; \tfrac{1}{1+\kappa} \,\Big\},
\end{equation}
the draw's own event \eqref{eq:draw} on the weight $Q$ alone, with no occurrence of $H$ in it. The whole $\kappa$-family of Accordance constraints collapses onto the family the draw already supplies: on an independent attribute the principle adds nothing, the two weightings agree, and the inference returns nothing. This is the framework's Theorem of Aimless Optimation---its control on scope, the guarantee that the method is silent where no observer-relevance is present---recovered as the principle's own value at independence rather than posited beside it, and available only because the principle carries $\kappa$: at one threshold a bound can hold with slack for a strongly relevant attribute, and it is the identity of the two events at every $\kappa$ at once that shows nothing about $H$ survives.

\textit{Tightness, and the realized member.} The bound is tight, attained as the pair's weights tend to $(1/(1+\kappa),\, \kappa/(1+\kappa))$; the event is non-strict for $\kappa > 1$, as the principle it descends from is, and strict at $\kappa = 1$, since on a two-cell partition an exact tie carries positive probability: at $q = \tfrac12$ it is certain. Note also that $A$ names the \emph{realized} member throughout, so both sides of the event are functions of the draw; the threshold is not fixed in advance. Display~\eqref{eq:accordance} is then the $\kappa = 1$ reading of \eqref{eq:exceedance} --- the realized member sits on the majority side of its pair except on posterior measure at most one half, the dropped point in the larger span --- and the content of the companion's own gloss that the principle ``holds more often than not'' \cite{Nexus} --- and $\lesssim$ is henceforth \emph{defined} by \eqref{eq:exceedance}: violations are not forbidden but improbable, with $\kappa$-fold exceedances confined to probability $1/(1+\kappa)$ and the entire tail profile thereby fixed, no free tolerance anywhere.

\textit{The canonical pair, and finer partitions.} The canonical pair is the attribute against its complement --- $P(A') = 1 - P(A)$, the form in which the companion boxes the principle \cite{Nexus} --- and for that pair the bound is tight. Finer-grained alternatives arise by refining the crossroads into a larger partition, where the mediocrity claim becomes that the realized cell stands near the partition's mode; the bookkeeping for refined partitions is the companion's, and every application in this paper runs on complement pairs.

% R864 (Opus 5.5, Oct 26 2026), at the author's direction (his message of October 26, 4:36 PM ET, and its follow-up):
% the chain's weights stated as the cosmic weighting's, exactly; the prior read as the provenance note's Earth count in
% exact form; the padding route closed; and the background state shown to leave the bounded event unchanged at a fixed
% target (the companion's Proposition on ambient invariance), so that the parent-conditioned and unconditioned forms are
% one principle.
% R866 (Opus 5.5, Oct 26 2026), at the author's direction: the background at a crossroads stated as a common root, after the companion's Common
% Root Principle (it read "the latest state $E$ at which $A$ is still blank"); likewise the companion's parent state in
% the paragraph below (it read "the latest state at which the attribute is still blank").
\textit{The weighting, and the background state.} Every probability in the display is the cosmic weighting's, exactly: one admissible cosmic measure, taken for the whole chain, and no measure tailored to the target or to the attribute, which is what cosmic-mediocrity forbids. The Aethae of \eqref{eq:chain} are stated content, weighted as that measure weighs them, so that $Q$ is the weight of $A$ among the states realizing the target, and $P(A)$ is the provenance note's Earth count in its exact form: the weight of the attribute among the Earths of the background state, every planet an Earth reached by perturbation (Section~\ref{sec:theorem}). The same fixity keeps a target from being built to suit. A target padded to lighten the realized cell gains nothing from content the cosmic weighting does not realize, since that content adds no weight; content it does realize, added after the realization in order to lighten that cell, makes a partition chosen by its outcome rather than the declared blank over which the draw is taken, and the draw says nothing of it: a selection after looking is covered only by counting its whole candidate family (Section~\ref{sec:multiplicity}), and a padding chosen freely has no family to count. The declared background state, by contrast, can be widened or narrowed without changing what is bounded, so long as the target stays inside it. For a target $H^{+}$ of the chain and any state $E$ containing it, $H^{+} \sqsubseteq E$, in which $A$ and $A'$ both carry weight,
\begin{equation}
\label{eq:background-invariance}
\frac{P(H^{+} \mid A', E)}{P(H^{+} \mid A, E)} \cdot \frac{P(A' \mid E)}{P(A \mid E)} \;=\; \frac{P(H^{+} \cap A')}{P(H^{+} \cap A)} \;=\; \frac{1-Q}{Q}, \qquad Q = P(A \mid H^{+}),
\end{equation}
the factors $P(A \cap E)$, $P(A' \cap E)$ and $P(E)$ cancelling. The canonical ratio and the prior odds change with $E$, by one common factor, and their ratio does not, so the event the draw bounds is fixed by the target alone. At a crossroads the background is a common root $E$ of its branches, an Aethus given by stated content of which both are children, at its narrowest their union, the provenance note's branch point; a moment of time, as in a history rewound and run again, names no Aethus and is no root \cite{Nexus}. The target is the observers arising from the root, $H^{+} = H \cap E$, a member of the chain because the reasoner's own state lies inside it; read at that target with every probability unconditioned, the principle bounds the same event, so the parent-conditioned form and the unconditioned form are one principle \cite{Nexus}. Moving the target is another matter, since the target fixes $Q$: within the chain the draw holds at every target, each with its own $Q$, and above $H$ it holds under the screening condition stated with the rarity bound (Section~\ref{sec:rarity-bound}).

One correspondence with the companion should be recorded, because the companion derives the same display by a differently grounded route and a reader of both documents should see one argument, not two. There the derivation proceeds in two legs \cite{Nexus}: the mediocrity step is grounded at the \emph{parent state} of each crossroads --- a common root of its branches, given by stated content --- as the companion's Copernican draw, on which the realized refinement is distributed by the parent-conditional weights; the anthropic leg then enters as the \emph{observer tilt} on that draw, the operative distribution being the parent draw reweighted by each branch's observer-productivity, and the boxed form falls out of the tilted draw's typicality requirement in two lines, every term parent-conditional. The present derivation is the self-contained composite: P1 as stated here takes the tilted, observer-conditioned draw as its premise directly, so the companion's two legs appear as one. The exceedance form is available identically on either grounding, and the two agree on the bound; the companion's parent-free reading, every probability unconditioned, needs no license of its own, since at a common target it bounds the event the parent-conditioned reading bounds, \eqref{eq:background-invariance}.
% R864 (Opus 5.5, Oct 26 2026): the parent-free reading's license removed, as at the source, where the lineage-screening
% license was retired with the old global form (R251); the clause read "the companion's global, parent-free reading
% additionally names its own license --- a lineage-screening assumption --- which the local form used here never needs."

One convention must be declared, because the self-sampling/self-indication fork re-enters through it if left silent. $H$ is graded --- the production of observers --- so $P(\cdot \mid H)$ needs a reading: the convention adopted here, matching the companion's usage, is the \emph{existence} reading, on which a branch's observer-productivity is the probability that observer-production occurs on it, not the expected observer count. A count-proportional reading would relocate self-indication's weighting to the state level and invite a state-level presumptuous philosopher, since an attribute multiplying observers would acquire enormous tilt, and attribute-locality does not prevent an attribute from being globally consequential. Two things follow and should be split explicitly: the exceedance bound is convention-invariant --- it requires only some fixed observer-conditioned weighting on the pair --- so the theorem below stands on either reading; what is convention-sensitive is the interpretation of the forced canonical ratios, which is link (ii)'s territory and is treated there.

% R845: the direction licensed by the model, not by Bayes' theorem (it read "directional in an unfamiliar way").
Three features of the principle do the work below, and each deserves one sentence. First, the constraint is \emph{local}: it binds attribute by attribute over a partial state, and at no point is a global optimization over complete histories performed or needed---there is no maximization step, only a family of pairwise bounds. Second, the constraint is \emph{directional} in a less familiar way: it licenses inferences from an attribute's prior improbability to its observer-relevance, the reverse of the selection-effect correction that reasons from observer-relevance to distorted frequencies. Bayes' theorem runs both ways; the model, not the theorem, licenses this reading. Third, the constraint is \emph{quantitative}: the bound in~\eqref{eq:accordance} has magnitudes on both sides, so the method it supports---taking a known feature of the record, asking what its realized-over-alternative prior ratio was, and reading off the forced observer-relevance ratio---produces bounded inferences rather than a binary verdict. Schematically: if the realized attribute was a hundredfold the a priori unlikelier member of its pair, the bound forces the alternative to have been, to the same order, the less observer-productive member---a magnitude, portable into further reasoning, where the bare selection-effect truism would have delivered only a shrug. The companion names this the method of counterfactual quantification and runs it at length \cite{Nexus}; here only the principle itself is needed.

\section{Two relocations}
\label{sec:relocations}

Because the vocabulary of typicality invites assimilation to observation-selection theory, the difference deserves explicit statement against the tradition by name. The anthropic principle of Carter \cite{carter1974large}, and the observation-selection apparatus systematized by Bostrom \cite{bostrom2002anthropic}, reason over reference classes of observers situated within complete histories; the tradition's central internal debate---self-sampling against self-indication, in Bostrom's terms \cite{bostrom2002anthropic}---is a debate about \emph{which} such class to sample from, with both parties holding the classical unit of reasoning fixed: a determinate timeline equipped, in effect, with a binary indicator for whether it houses observers.

The completed unit is not an idle assumption; it is where the tradition's characteristic troubles are manufactured, on both horns of its central debate. Self-sampling over a completed history is what exposes the Doomsday inference in its most criticized form---one's birth-rank treated as a random draw from a completed roster \cite{leslie1996end, bostrom2002anthropic}---while self-indication, adopted largely to escape it, purchases the escape at the price of the presumptuous-philosopher pattern, weighting hypotheses by the observer-counts of completed worlds \cite{bostrom2002anthropic}. Each horn inherits its embarrassment from the same source: a reference class read off a totality.

The present framework differs by two relocations. First, the reference class ranges over observer-situations \emph{across possible histories}, not over positions within one fixed history. That first relocation is, by itself, continuity rather than departure --- Bostrom's own fine-tuning applications run the Self-Sampling Assumption across possible worlds \cite{bostrom2002anthropic} --- and the provenance note above says the same from the biographical side: the principle was first derived on a frankly cosmological Earth-count picture, squarely inside the tradition's furniture. The departure is the second relocation. Second, the object of reasoning is the partial state itself, never its completion. From this vantage the standard apparatus is the artificially completed limiting case, and the tradition's most persistent difficulty---the reference-class problem---is, on the reading this paper proposes, an artifact of the completion: a debate over which classical class to sample from cannot close, because the classical class was never the right unit. The diagnosis extends to the tradition's characteristic rescue moves. A reference class welded to complete histories, when its explananda outgrow one history, has a single direction available: manufacture more complete histories and reason over the enlarged stock. Such constructions are epicycles in the strict historical sense, machinery added to preserve a background commitment rather than machinery demanded by the phenomena; the present move is not a further epicycle but the replacement of the deferent.

The canonical formalization of self-locating typicality must be engaged by name before the calibrations run: Elga's indifference principle --- equal credence over subjectively indistinguishable centered worlds, the lesson of Dr.~Evil \cite{elga2004defeating}. The Earth-count picture recorded in the provenance note is recognizably an indifference intuition of Elga's kind, reached without contact with the formal literature --- though not in a vacuum: the seed was popular exposure to the Doomsday argument alone, and the direction taken diverged from the tradition's precisely because the reaction to that much was phenomenological rather than textual. The mature principle's differentia is then this. The Copernican draw is \emph{attribute-wise}, executed over a partial state's declared alternatives, and it never requires --- or permits --- the total centered-world partition Elga's principle quantifies over. Where that partition cannot be formed, indifference over it is not weakened but, under the admissibility principle below, \emph{unformulable}. The question itself stays askable: a hypothesis about how many observers there will ever be can be stated and given a prior. But the indifference Elga's principle needs ranges over the reasoner's place in that total, and a partial state that leaves the total open has declared no alternatives for it to range over, the total being blank relative to the state rather than divided among weighted alternatives. That last step is a principle, not a consequence of partiality alone. A partial state can leave the total blank while carrying alternatives for it---a prior over eventual population sizes---and alternatives for an event involving it, such as one's rank within it, and nothing in partiality forbids an indifference over the latter. What excludes it is an \emph{admissibility principle}: the draw ranges only over attributes whose weights the state's cosmic weighting fixes without passing through a partition of completed rosters. The criterion cannot be the answer space. The proposition that one's rank lies in the first tenth of the eventual population has the fixed answer space yes or no even while the population is unknown; its truth depends on the total, its answer space does not. What it lacks is a weight the state supplies: its two answers are weighted only by summing over ranks within completed totals, through the centered partition the draw does not permit, and the indifference over rank that sum requires, one's birth rank read as a uniform draw from a completed roster, is the step of the self-sampling assumption the Doomsday argument needs and the framework does not posit. The framework's own typicality is self-locating in a different sense: one's observer-situation is treated as typical among observer-situations across possible histories, conditioned on a history's producing an observer of the kind in question rather than weighted by how many it produces, and reasoned on the partial state attribute by attribute; a position in a completed roster has no attribute-wise counterpart, so the draw supplies no rank likelihood and no Doomsday update. The two commitments are these. \emph{Aetherian sampling}: of the category of mediocre consciousnesses taken abstractly, the Aetherian category, across the entire Open Cosmos, you are a random sample. \emph{Aetherian conditioning}: the physical universe conditioned on an observer-situation of that category resolves, step by step, as the observer-production Nexus itself resolves, so that observer-production $H$ may serve as the conditioning event throughout without a separate transport assumption at each link; this paper uses the companion framework's account of that step-by-step resolution \cite{Aethus} only through this principle, and asserts nothing further about it. Both are assumptions, and the use of $H$ throughout stands or falls with the second. The principle is argued here, not derived, and the Doomsday verdict below is stated conditional on it. The headline exhibit below is that fact run on the tradition's most famous argument.

% R845: Neal's full non-indexical conditioning; Hartle and Srednicki on typicality as a stated prior and on joint
% testing; what is distinctive placed in the model and procedure.
Two positions come nearer this paper's aims. Neal's full non-indexical conditioning conditions on all of one's evidence rather than on a reference class, usually returning what self-sampling and self-indication give together with a broad class, and, like the admissibility principle here, yields no Doomsday shift \cite{Neal2006}. Hartle and Srednicki deny that typicality may be assumed without evidence of a physical selection process, holding it a prior to be stated \cite{hartle2007typical}, and test a theory together with its distribution over our location \cite{srednicki2010large}. This paper states its typicality premise, (P1), and tests a conjunction including it; what is distinctive is not the conditioning, which is ordinary, but the attribute-wise formulation, the admissibility rules, the sampling commitments and the applications.

Three calibrations locate the framework among its neighbors. It is consistent with the weak anthropic principle as a selection-effect truism \cite{bostrom2002anthropic}; what it adds is the directional, quantitative constraint of Section~\ref{sec:principle}, which the truism does not supply. And its relation to the Doomsday argument---in the Carter--Leslie form \cite{leslie1996end}, and in Gott's Copernican variant \cite{gott1993implications}---is generalization of the method with localization of the claim: attribute-wise bounds over partial states, in place of a single global position-in-duration inference. 

\paragraph{The headline exhibit: birth rank, and the Doomsday argument's non-formulability.} Partiality's clearest bite, under the admissibility principle, deserves top billing rather than a calibration's corner, and it is this. The global birth-rank inference is itself not recovered: a position in a completed roster has no attribute-wise translation --- and the retort that ``being in the first ten percent of all humans'' could itself be made an attribute fails on inspection, since, although its answer space is a fixed yes or no, its weights are fixed only by summing over ranks within completed rosters, the partition it was meant to replace --- and the framework's verdict on its founding relative is adoption of the method with retirement of that application's unit. The tradition's own nearest approach to the present method is Carter's hard-steps model \cite{carter1983anthropic}: quantitative reasoning from a typicality-flavored premise --- the coincidence of evolutionary and stellar timescales --- to rarity-flavored conclusions about the difficulty of critical evolutionary transitions, the analysis on which the filter framing builds \cite{hanson1998greatfilter}. The relation is generalization with relocation: where the hard-steps model reasons over durations within one timeline, with the steps as its posited unit, the Accordance constraint binds attribute by attribute across histories --- recovering hard-steps-style conclusions where the record's improbabilities are duration-linked, and reaching attributes the duration argument cannot see. The modern Bayesian hard-steps literature is this application's live neighbor: the transition-timing analysis of Snyder-Beattie, Sandberg, Drexler, and Bonsall \cite{snyderbeattie2021timing}, together with the recent critical reassessment of Carter's model by Mills, Macalady, Frank and Wright, who argue that the timing of the supposedly hard transitions may reflect the sequential opening of global environmental windows of habitability rather than intrinsic improbability \cite{mills2025hardsteps}. That reassessment bears on this paper directly and should be met rather than noted: it challenges the empirical antecedent, that a natural history supplies a long list of independently improbable attributes at estimable prior, by proposing that the apparent improbabilities are artifacts of when the environment permitted each step. On this paper's reading the dispute relocates into the attribute vocabulary and the couplings: a step that was ``waiting for its window'' is an attribute whose cosmic prior is high once the window is stated and whose observer-relevance is carried by the window---which is not a refutation of the method but an instance of it, and one the comprehensive development's coupled-structure reading takes up \cite{Nexus}. The attribute-wise draw enters that debate with a definite posture: step-hardness claims are couplings over declared alternatives, evaluated per stock, so the reassessment controversy is, from here, a dispute about the attribute vocabulary --- which is where the framework says such disputes live.

Indeed the framework's own provenance runs through the Doomsday argument and nothing else in the tradition, a point recorded in the provenance note below.\footnote{A placement aside, for the reader tracking the wider program: this paper is the fifth face of a single two-winged framework, and the first drawn from its anthropic wing---the four prior extractions, all from the Aethic wing, treat a counterfactual semantics, an informational classification of the quantum--classical boundary, the composition primitive that classification turns on, and a decision rule; they are surveyed together, connections drawn, in the companions \cite{Aethus, AethusBrief}. The two wings are logically independent: the overview \cite{AethusBrief} makes no use of the anthropic apparatus, and where a sibling extraction touches it --- the decision-theoretic face's second route --- it does so as an explicitly non-presupposed alternative; the wing's own development is \cite{Nexus}, and the independence is structural, as the provenance note records.}

\section{From the principle to a bound on observer-abundance}
\label{sec:rarity-bound}

Two claims are in play here, and they are distinct. The first is structural: the Accordance Principle bounds the probability that a biosphere reaches observers by the prior of any attribute realized in one's own record, at a stated confidence. That is a theorem, and it is proved in four lines. The second is empirical: that Earth's record supplies attributes whose priors are small enough for the bound to bite. That is argued below, from the structure of a natural history rather than from any one of its dramatic events.

The principle \eqref{eq:exceedance} says that the event $P(A')\,P(H \mid A') \ge \kappa\, P(A)\, P(H \mid A)$ has probability at most $1/(1+\kappa)$, so its reverse holds with probability at least $\kappa/(1+\kappa)$. Since $A$ and $A'$ partition the space, the law of total probability gives
\begin{equation}
\label{eq:total}
P(H) \;=\; P(A)\,P(H \mid A) + P(A')\,P(H \mid A'),
\end{equation}
and on the event of probability at least $\kappa/(1+\kappa)$,
\begin{equation}
\label{eq:bound-steps}
P(H) \;<\; (1+\kappa)\, P(A)\, P(H \mid A) \;\le\; (1+\kappa)\, P(A),
\end{equation}
the last step because a conditional probability is at most one.
The step just taken has two parts of very different weight, and separating them locates where the anthropic content sits. The inequality itself is elementary: since $P(A) \ge P(A \cap H) = P(A \mid H)\,P(H)$,
\begin{equation}
\label{eq:elementary}
P(H) \;\le\; \frac{P(A)}{P(A \mid H)} \;=\; \frac{P(A)}{Q},
\end{equation}
which holds for any $A$ and $H$ whatever, is a rearrangement of the definition of conditional probability, and asserts nothing anthropic. What the framework supplies is not this inequality but the \emph{calibration of its denominator}: the mediocrity principle says that the realized $Q$ is super-uniform in its lower tail over the draw, which is what converts \eqref{eq:elementary} from an identity into a bound with a stated confidence. The Accordance Principle is then what connects that calibrated denominator to observer-relevance. With that division in view:
\begin{theorem}[The rarity bound]
\label{thm:rarity}
Let $A$ be an attribute realized in the reasoner's record, $Q = P(A \mid H)$, and suppose $Q$ is super-uniform over the draw in the lower tail, $\Pr[Q \le w] \le w$ for $w < \tfrac12$. Then for every $\kappa > 1$, with probability at least $\kappa/(1+\kappa)$,
\begin{equation}
\label{eq:rarity}
P(H) \;\le\; (1+\kappa)\, P(A).
\end{equation}
The inequality \eqref{eq:elementary} carries no assumptions; the confidence attached to it is the lower-tail super-uniformity, and that is the theorem's whole anthropic content. In coverage terms: $P(H)$ is a fixed feature of the cosmic weighting, what is random across realizations is the bound $(1+\kappa)P(A)$ the realized attribute produces, and the statement is that this bound is a $\kappa/(1+\kappa)$ upper confidence bound on $P(H)$---the realized record generates the bound, and the bound covers the fixed quantity with the stated probability.
\end{theorem}

\begin{lemma}[Fixed parameter]\label{lem:fixed}
The tail of Theorem~\ref{thm:rarity} is uniform in every \emph{predetermined} $\kappa > 1$: the event's threshold $\kappa\,P(A)/P(A')$ varies with the realization because the statistic varies, and that is licensed. What is not licensed is $\kappa = \kappa(A)$, a threshold parameter chosen from the realized member, which converts a bound proved for each fixed $\kappa$ into a bound at a random $\kappa$ and requires a separate simultaneous-bound argument over the values $\kappa$ can take---the union bound of Section~\ref{sec:multiplicity} being the available one. The withdrawn prior bound was the instance $\kappa = (1-P(A))/P(A)$.
\end{lemma}

% R845: the screening license for targets above H (it read "The same bound holds at any conditioning target H+ of the
% chain in place of H, so it bounds the abundance of water-bearing worlds").
\noindent At $\kappa = 1$ the bound is $P(H) \lesssim 2\,P(A)$ with probability at least one half; at $\kappa = 9$ it is $\Pr[P(H) > 10\,P(A)] \le 0.1$. The multiplier and the confidence rise together, the trade being the quantile ladder \eqref{eq:ladder} read on $P(H)$. In words: if the realized attribute were much rarer than the observer-abundance it sits inside, the observer-mass would have had to lie on the alternative, and mediocrity would have placed the reasoner there. The same bound holds with an ancestor $H^{+}$ of $H$ in place of $H$ under the screening condition Section~\ref{sec:principle} calls for, observer-production bearing on $A$ only through $H^{+}$; under it the bound limits the abundance of water-bearing worlds by an attribute upstream of water as readily as that of observers, and an attribute favorable to observers through a channel bypassing $H^{+}$ breaks it \cite{Nexus}. Read in a background state $E$, at the target $H \cap E$ of the observers arising from it, the same four lines give, by \eqref{eq:background-invariance} and at the same confidence, $P(H \mid E) \le (1+\kappa)\,P(A \mid E)$ for any attribute $A$ open at $E$ and realized in the record: a bound on the step from the earlier stage $E$ to observers, by the weight among $E$-states of an attribute realized between them. The ancestors include generalizations of observer-production itself, stated as Aethic content rather than as lineages, such as mammal-grade, multicellular or eukaryotic life, each with observer-production within it as one of its cells, $H \equiv_{\mathcal{V}} A_{H}\,H^{+}$, where $A_{H}$ is the attribute of going on to produce observers and the two cells are weighted by the cosmic weighting exactly. The screening condition is then the cosmic independence of $A$ and $A_{H}$ within $H^{+}$, $P(A \mid H) = P(A \mid H^{+})$, under which the realized cell's weight at $H^{+}$ is its weight at $H$; read in a background state $E$, at the target $H^{+} \cap E$, the condition is the same independence within $H^{+} \cap E$, which neither implies nor is implied by the cosmic one. Down a ladder of such grains, nested by leaving further content blank at each, each coarser grain the background of the next finer and each step under the screening condition at its own target, the bound prices every step and not only the last, the $n$ bounds of a predeclared ladder holding together except with probability at most $n/(1+\kappa)$ \cite{Nexus}: with $A$ the planet's bearing liquid water, the background all Earths and the target multicellular life, it bounds the frequency of Earths that reach multicellular life by a multiple of the frequency of those that bear water.
% R872 (Opus 5.5, Oct 27 2026), at the author's direction (the Pr[] form the default): the rung at kappa = 9 written
% in the exceedance form; it read "at $\kappa = 9$ it is $P(H) \le 10\,P(A)$ with probability at least $0.9$" (R871's
% subscripted forms withdrawn).
% R864 (Opus 5.5, Oct 26 2026), at the author's direction: the bound read in a background state, where it prices a
% step (the companion's step bound), and the ancestors above observer-production, H = A_H H^+ with the cells weighted
% by the cosmic weighting exactly, the screening condition read as cosmic independence.

\textit{An instance, sourced.} Kaiho and Oshima estimate that a Chicxulub-class impactor could have produced an extinction of the end-Cretaceous severity only by striking the roughly thirteen percent of Earth's surface underlain by hydrocarbon-rich sedimentary rock, whose burning supplies the soot that drives the cooling \cite{Kaiho2017}. Take the realized attribute to be the impact's having struck such a site, with cosmic prior $P(A) \approx 0.13$. The figure is their model-dependent estimate for an impact on Earth under their stated conditions; reading it as a cosmic prior is this illustration's assumption, which a population argument would have to underwrite, and not a consequence of their estimate. Substituting, $P(H) \lesssim 0.26$ at $\kappa = 1$ with probability at least one half, and $\Pr[P(H) > 0.52] \le 0.25$ at $\kappa = 3$. One published figure, one attribute, one bound at a confidence one chooses: that is the procedure, and the figure's own uncertainty is carried into the bound rather than hidden by it.
% R871 (Opus 5.5, Oct 27 2026), at the author's direction: the prior, the source's "roughly thirteen percent", written
% "\approx" (it read "\lesssim 0.13", the tendency symbol being reserved for priced tendencies).
% R872 (Opus 5.5, Oct 27 2026), at the author's direction (the Pr[] form the default): the bound at kappa = 3 written
% in the exceedance form (it read "$P(H) \le 0.52$ at $\kappa = 3$ with probability at least $0.75$"; R871 had the
% subscripted symbol).

\textit{Many attributes.} A single attribute bounds $P(H)$ by a small multiple of its own prior, and no single attribute in a natural history has a prior small enough to make that bound the Rare Earth hypothesis. What does is a conjunction, used through the predeclared-event test of Section~\ref{sec:test}. Let $\mathcal{A}_1, \dots, \mathcal{A}_m$ be binary attribute partitions fixed before the realization draw, with rare sides $A_i$ of prior $p_i = P(A_i)$, and let $C = A_1 \cap \cdots \cap A_m$ be the event that every rare side is realized. The test needs one assumption about these events, an upper bound on the prior of their conjunction:
\begin{description}[leftmargin=3.2em, style=nextline, itemsep=0.2em]
\item[(A1) Cosmic submultiplicativity.] $P(C) \le \prod_i p_i$, under the weighting in which nothing has been conditioned on; independence of the attributes under that weighting is the equality case.
\end{description}
\noindent It needs nothing about the observer-conditioned proportions. The predeclared conjunction $C$ is a fixed event and $P(C \mid H)$ a fixed number; what is random is which member of $\{C, C^{c}\}$ the record realizes, and the test is a statement about that. Where a conjunction is instead read off the record, the realized pattern is a cell of a partition with up to $2^m$ cells and does not inherit a binary lower tail: for two independent fair coins each conjunct's proportion is $\tfrac12$, yet the realized pattern has proportion $\tfrac14$ on every draw. That case is priced by the realized-cell bound, Lemma~\ref{lem:realized-cell}, at the cost of the number of cells.



\subsection{The test}\label{sec:test}

\noindent The paper's main result is the test of Corollary~\ref{cor:rareearth} below. It needs only {\rm (A1)}.

\begin{lemma}[Realized-cell bound]\label{lem:realized-cell}
Let $\mathcal{P}$ be a predeclared partition with at most $K$ cells and $C(X)$ its realized cell. Then $\Pr\big[P(C(X) \mid H) \le \delta/K\big] \le \delta$, so that with coverage at least $1-\delta$,
\begin{equation}\label{eq:realized-cell}
P(H) \;\le\; \frac{K}{\delta}\, P\big(C(X)\big).
\end{equation}
For $m$ predeclared binary attributes $K \le 2^m$; and when the realized cell is the all-rare one, or any cell whose prior is at most $p^m$, {\rm (A1)} gives $P(H) \le (2p)^m/\delta$. A realization with common sides has a larger prior, and its bound is \eqref{eq:realized-cell} with that prior.
\end{lemma}
\begin{proof}
Under the draw, $\Pr[P(C(X) \mid H) \le \delta/K]$ is the $H$-conditioned mass of the cells whose conditional probability is at most $\delta/K$, at most $K \cdot \delta/K = \delta$; and $P(H) \le P(C)/P(C \mid H)$ by \eqref{eq:elementary}.
\end{proof}
\noindent The lemma needs no assumption about the observer-conditioned weighting, and it makes the price of selecting a realized pattern explicit as the factor $K$; its threshold is $p < \tfrac12$. For a predeclared conjunction the test is simpler still: under the null $P(H) \ge h$, $P(C \mid H) \le P(C)/h$, so rejecting the null when $C$ occurs errs with probability at most $\delta$ whenever $P(C) \le \delta h$. What neither result licenses should be stated as plainly. The event that every member of an indefinitely extended family realizes its rare side has, under any model with $P(H) > 0$, observer-conditioned probability $P\big(\bigcap_m D_m \mid H\big) = 0$, since $P(D_m \mid H) \le p^m/P(H)$; conditioning on such a record does not retain the finite coverage interpretation. The defensible statement is finite: each predeclared family carries a specified error rate, and a particular record can yield a small bound. Inspecting further attributes until a bound crosses a target is a stopping rule, and fixed-$m$ coverage does not by itself give coverage simultaneous over $m$.

\noindent \textit{A worked case.} Take $m = 30$ predeclared attributes, each with rare-side cosmic prior $p = 0.2$, and a record realizing every rare side. The conjunction rule of Corollary~\ref{cor:rareearth} gives $P(H) \le 0.2^{30}/\delta$: about $1.1 \times 10^{-20}$ at $90\%$ and $1.1 \times 10^{-19}$ at $99\%$; at $p = 0.1$, about $10^{-29}$ and $10^{-28}$. The realized-cell rule (Lemma~\ref{lem:realized-cell}), which would also have rejected on any other cell of low enough prior, pays $2^{30}$ for that reach: $P(H) \le 0.4^{30}/\delta$, about $1.2 \times 10^{-11}$ at $90\%$ and $1.2 \times 10^{-10}$ at $99\%$, and about $1.1 \times 10^{-20}$ and $1.1 \times 10^{-19}$ at $p = 0.1$. Dropping any attribute whose prior, or whose place under {\rm (A1)}, a reader doubts costs one factor of $p$ under the conjunction rule and of $2p$ under the realized-cell rule, which is the separability both routes were chosen for.

\begin{table}[ht]
\centering
\small
\caption{The candidate inputs of the test, one row per attribute, with the figures the published literature supplies as of September 25, 2026 and their status. A ceiling is an upper bound on the rare side's cosmic prior; a larger value keeps the test valid and weakens it. A figure whose status is a central estimate or a model input is a candidate, not a ceiling, and the row does not certify the test's conditions. The third column states what {\rm (A1)} would require of the row, and the last records temporal precedence only, which is weaker than the corollary's condition (see the text). The oxygenation row has no non-anthropic figure and is not used; the last row is the worked case, illustrative. The companion long paper collects the sources and the correlations \cite{Nexus}.}
\label{tab:application}
\begin{tabular}{p{0.26\textwidth} p{0.17\textwidth} p{0.27\textwidth} p{0.2\textwidth}}
\hline
Attribute (rare side, joint form) & Figure, and its status & What {\rm (A1)} would require & Identified before the rarity argument by \\
\hline
A Chicxulub-class impact occurs and strikes hydrocarbon-rich rock at a steep angle, within the window of complex life & per impact $\le 0.065$ (site $0.13$ \cite{Kaiho2017} $\times$ angle $\le 0.5$ \cite{collins2020chicxulub}); per history about $0.5$ at the conservative flux \cite{nesvorny2021dark}; model estimates, the per-history value assuming that flux and a $1.5$~Gyr window & a ceiling on the joint conditional on the architecture row, which is unsupplied: the impact flux may depend on the giant planets & the kill mechanism, named before any rarity argument; the site and angle themselves describe Earth's realized impact, so the selection rule is open \\
A cold Jupiter on a wide orbit and no super-Earth or sub-Neptune inside Mercury's orbit & about $0.10$ \cite{mulders2018epos, zhu2018thirty, fulton2021cls, bryan2019excess, rosenthal2022cls3, bonomo2023cold}; occurrence-rate estimates with completeness corrections, not upper bounds & the joint of the two rare sides lies below their product only if their negative correlation holds under the cosmic measure; priced as one attribute & giant-planet dynamics and inner-system clearing \\
A moon massive enough to stabilize obliquity & $1/4$, the high end of a simulation range running from $1/45$ to $1/4$ with a central value above $1/12$ \cite{elser2011moon, laskar1993stabilization}; a model input, not a bound on the cosmic prior & a ceiling conditional on the architecture row; the correlation is unknown, so none is supplied & obliquity stabilization, weakened by \cite{lissauer2012obliquity} \\
Oxygenation early enough relative to the star's brightening & none; assumed distributions give $0.4$--$0.85$ \cite{catling2005oxygen, bixel2020oxygen} & --- & the oxygen requirement of animals; no cosmic prior yet \\
Worked case: thirty attributes & $0.2$ each & assumed & $m$ fixed by $\varepsilon$ and $\delta$ in advance \\
\hline
\end{tabular}
\end{table}

\begin{table}[ht]
\centering
\small
\caption{What each candidate input must supply before it is a ceiling the corollary can consume, and where each stands. ``Open'' marks a task, not a certification.}
\label{tab:status}
\begin{tabular}{p{0.13\textwidth} p{0.17\textwidth} p{0.15\textwidth} p{0.19\textwidth} p{0.13\textwidth} p{0.12\textwidth}}
\hline
Factor & Event and population & Source estimand & Bound status & Conditional on prior rows & Selection rule; transport \\
\hline
Impact (site, angle) & a lethal-by-site-and-angle impact within the window of complex life, per history & surface fraction and arrival geometry, model results; per-history value at a stated flux and window & estimates; the window and flux are assumptions & on architecture: open & outcome-selected from Earth's impact; family and correction open; Earth models to the cosmic measure: open \\
Architecture & cold Jupiter and clear inner system, per FGK star & occurrence rates & estimates with completeness corrections & first row & survey-defined; nearby FGK stars to the cosmic measure: open \\
Moon & a stabilizing moon, per terrestrial planet & simulation frequency & $1/4$ is the high end of $1/45$--$1/4$ & on architecture: unknown & categorical; simulated systems to the cosmic measure: open \\
Single star & single host, per star & multiplicity fraction & survey estimate of $P(\text{single} \mid \text{star})$, not of $P(\text{single} \mid \text{biosphere})$ & unknown & categorical; solar-type stars to biospheres: open \\
Rocky HZ planet & at least one rocky planet in the conservative zone, per star & planets per star, $0.37$--$0.60$ central, $68\%$ intervals to $1.5$ & central estimate; $P(N \ge 1) \le \mathbb{E}N$ gives no ceiling below one; and a biosphere presupposes the planet & on single star: open & categorical; Kepler FGK stars to the cosmic measure: open \\
\hline
\end{tabular}
\end{table}

\noindent\textit{Predeclaration, and what the last column shows.} The corollary's condition is that the attributes be fixed independently of the realization, not merely before a rarity argument about it. The table's last column records the weaker fact, temporal precedence: the surveys counted stars, moons and architectures for their own reasons, and the kill mechanism was named before anyone priced it. An attribute can carry an excellent physical motivation and still have been chosen because it describes Earth's known history, and studying the mechanism of a known impact does not by itself yield an outcome-independent rule for choosing that impact's attributes. Three routes meet the condition: a selection rule demonstrably independent of the outcomes tested; a specified larger family of candidate attributes with a valid selection correction, which is what Lemma~\ref{lem:realized-cell} prices as the factor $K$; or a calibrated procedure that incorporates how the attributes, thresholds and windows were chosen. The impact row fails the condition as it stands, its site and angle being read off the realized impact; its family is specifiable, the attributes the habitat-model literature considers for a Chicxulub-class impact, size, angle, site lithology, season and coincidence with a volcanic province, and the correction is the realized-cell one, but until the family is written down and priced the row is a candidate. The categorical rows come closer, though ``Solar-System-like'' as a survey target is itself outcome-shaped. Reporting common features beside rare ones, as the catalog does, is honest bookkeeping and establishes none of the three. A fixed named conjunction has a fixed probability, and the realized-cell calibration transfers to it only under a complete reporting rule that states what is reported when the conjunction fails; joint error control over the draw is not conditional confidence after the rare outcome has been selected.

% R845: every scenario with its Poisson mean, as impact_example.py computes it; at the ratified inputs the realized-cell
% bound is 0.35 (it read "by nothing", a superseded row's value, and gave the superseded conservative 0.72 as the
% conservative inputs'); "defaults awaiting ratification" dropped, the inputs having been ratified on 2026-09-29.
\noindent\textit{The impact row, fully specified.} Run as one complete event on declared inputs, Chicxulub-class impacts arriving as a Poisson process of mean $N$ and each realizing the declared family independently (target rock $0.13$, an angle bin of $45$--$60^{\circ}$ at $0.25$, flood-basalt coincidence $0.042$, so $q = 0.001365$ per impact), the chain has per-history probability $1 - e^{-Nq}$. At the ratified inputs, the Phanerozoic window at the literature flux of $3$ impacts of $D > 10$~km per Gyr \cite{nesvorny2021dark}, both fixed on physical grounds, where $N = 1.62$, the probability is $0.0022$, and the chain bounds $P(H)$ by $0.044$ under the conjunction rule, if the family counts as predeclared, and by $0.35$ under the realized-cell rule at $K = 8$; at the conservative inputs ($N = 2.16$, site $0.20$, coincidence $0.049$), by $0.11$ and $0.85$; at ten impacts in $1.5$~Gyr ($N = 10$), by $0.27$ and nothing. The companion gives the computation \cite{Nexus}. One honest chain buys a factor of a few, and the window decides more than any attribute's read.

\noindent\textit{A conditional benchmark.} If the three sourced rows, a single host star at $0.56$ \cite{raghavan2010multiplicity} and a rocky habitable-zone planet at $0.6$ \cite{bryson2021eta} were ceilings on the conditional probabilities the chain rule needs, under one per-history measure, and the attributes had been selected independently of the outcome, the conjunction rule would give a product of about $4.2 \times 10^{-3}$ and, at $\delta = 0.05$, $P(H) \le 0.084$. None of the five presently meets that standard, and Table~\ref{tab:status} records where each falls short. The rocky-planet figure is the upper of two central estimates of planets per star, $0.37$ and $0.60$ under differing completeness assumptions, with $68\%$ credible intervals reaching $1.5$; the estimand bounds the probability of at least one such planet only through $P(N \ge 1) \le \mathbb{E}N$, which supplies no ceiling below one, and a biosphere presupposes the planet in any case, so the factor cannot multiply into a probability over biospheres. The moon's $1/4$ is the high end of a simulation range. The single-star fraction is a per-star marginal, not $P(\text{single} \mid \text{biosphere})$. The architecture and impact figures are occurrence estimates. The conditional ceilings that (A1) or the chain rule requires are unsupplied, and the rows sit in different populations, per impact, per history, per planet and per star, whose embedding in one measure over biospheres has not been shown. Two further points hold whatever the inputs: $\delta$ covers the observer draw only, not the uncertainty of the inputs, which needs its own simultaneous validity; and Bayesian credible intervals are not frequentist coverage guarantees. The figure is therefore a conditional benchmark of the rule's reach at inputs of this size, not an established bound, and the companion's further figures inherit every condition above. Correlated attributes, metallicity with giant planets and zone membership, stellar type with tidal locking and flare rates, would be handled by the chain rule once conditional ceilings were supplied, the survey rates being conditional on nearby FGK stars already inside the zone, which fixes the population for those rows and for no others. The end-Cretaceous row shows the distinction the test lives on: the per-impact narrowness, $3$--$7\%$, is the striking number, but a history can hold several such impacts, so the attribute the test consumes is the per-history one unless the timing joint, that the lethal impact fell when its successors were ready, is priced. Which of these attributes count as predeclared, given that Earth's record was read long ago, is the question the fourth column answers row by row: each was named as a bottleneck for a reason independent of its rarity, and the attributes that came out common, the impact's season at a prior of one quarter for instance, are reported beside them.

% R845: the 0.04 ceiling takes the benchmark to about 0.007, and the complete event's 0.0022 to about 4 x 10^-4 (it read
% "falls from about 0.08 to about 4 x 10^-4" under the first).
\noindent\textit{Where the program stands.} The benchmark above, conditional as Table~\ref{tab:status} shows, says less about the world than about the questions that have been asked of it. The literature prices categorical attributes, the kinds of star, moon and architecture the surveys were built to count; the contingent reads, the impact's site and angle, exist because a few groups asked a question of this paper's shape for their own reasons. The joint forms, the per-history prior and the sparing window are questions no one has been asking, so the present ceiling is set by practice, not by the record, which admits far more. Coupling inference cuts both ways: coupled attributes merge, so an honest $m$ can shrink, while a cascade can price an attribute the literature never touched and the companion's landing-depth route recovers the effective number of independent attributes from depths rather than from a count. Under the conjunction rule at $95\%$ confidence, a dozen to sixteen attributes at ceilings like the present ones would put $P(H)$ below one in a million, once they were ceilings on conditional probabilities under one measure and predeclared, and about $1.8$ times as many at $p = 0.2$ on the realized-cell route, which charges $2p$ per attribute of a declared family; the thirty of the worked case show scale, not need. The companion collects the computations that would supply them \cite{Nexus}. The first of them, run on published ages alone, finds that a random impact falls within $\pm 100$~kyr of a volcanic province's main eruptive phase about $4\%$ of the time, so that the realized conjunction, an impact lethal by site and angle falling within such a phase, carries a per-history ceiling of about $0.04$ against $0.5$ without it. If the coincidence counts as part of the rare side, the province's warming having mitigated the impact winter, and if neither event triggered the other, the benchmark above falls from about $0.08$ to about $0.007$, and to about $4 \times 10^{-4}$ with the complete event's per-history probability at the ratified inputs, $0.0022$, in place of the $0.5$.


\begin{corollary}[Rare Earth]
\label{cor:rareearth}
Let $\mathcal{A}_1, \ldots, \mathcal{A}_m$ be binary attribute partitions \emph{predeclared}---fixed before the realization draw, not merely before the analyst computes from the observed record---each with rare-side prior at most $p$ and jointly satisfying {\rm (A1)}, and let $\varepsilon > 0$ and $0 < \delta < 1$ satisfy $p^m \le \delta\varepsilon$. If the record realizes the rare side of every $\mathcal{A}_i$, then the hypothesis $P(H) \ge \varepsilon$ is rejected at level $\delta$: the test that rejects it on the predesignated event $C = \bigcap_i A_i$, every rare side realized, errs with probability at most $\delta$ whenever the hypothesis is true. Since a rare side's prior is below one---for a binary partition it is at most one half---every $\varepsilon$ and every $\delta$ are met by a finite $m$. The probability that a biosphere reaches observers is thereby bounded below any positive level, at any confidence short of certainty, by a long enough predeclared record realizing its rare sides: the Rare Earth hypothesis, from the Accordance Principle, the law of total probability, and the record, with no High-Contrast Earth or landscape-contrast hypothesis assumed and no independence assumption beyond {\rm (A1)}. The order of quantifiers is the corollary's content. The partitions are fixed first, and with them the rejection event; which side each realizes is then read off the record. Two rejection rules are available and are distinct. The \emph{conjunction rule} above rejects only on the predesignated event $C$. The \emph{realized-cell rule} of Lemma~\ref{lem:realized-cell} rejects on whatever cell the family realizes when that cell's prior is at most $\delta\varepsilon/2^m$, the common-side realizations contributing their large priors and the rare-side realizations their small ones; it pays the factor $2^m$ for that selection among cells, and on an all-rare record it needs $(2p)^m \le \delta\varepsilon$, hence $p < \tfrac12$. What the corollary does not license is the converse order---scanning the record for attributes whose realized side happens to be rare and declaring those the family---because ``the realized member had prior $p$'' is an outcome-dependent fact, and under the null case of \eqref{eq:null-case} a rare-side realization is the event on which the tail can fail; any post-selection from a larger scanned family is covered only by a simultaneous correction over that full family, Section~\ref{sec:multiplicity}. Nor does it license inspecting further attributes until the bound crosses a target, which is a stopping rule: $m$ is fixed by $\varepsilon$ and $\delta$ before the record is read.
\end{corollary}
\begin{proof} Under the hypothesis $P(H) \ge \varepsilon$, $P(C \mid H) = P(C \cap H)/P(H) \le P(C)/\varepsilon \le p^m/\varepsilon \le \delta$, the middle step by {\rm (A1)}: the test rejects on $C$, and $C$ occurs, given $H$, with probability at most $\delta$. The realized-cell rule is Lemma~\ref{lem:realized-cell} with $K = 2^m$, and on an all-rare record $p^m \le \delta\varepsilon/2^m$ is $(2p)^m \le \delta\varepsilon$. \end{proof}
\noindent The corollary and the bound it rests on now have one standing: each takes a predeclared list of attributes and a confidence and returns a number, and every application in this paper and its companion is of that kind. The conjunction rule needs only {\rm (A1)} and a prior ceiling below one; the threshold $p < \tfrac12$ belongs to the realized-cell rule, whose extra reach, rejection on records that realize some common sides, is what the factor $2^m$ buys. What the corollary asks of a record is finite: $m$ predeclared contingencies realized on their rare sides, with $m$ set by the target. What it does not ask is that a record answer every question on its rare side. Under any model with $P(H) > 0$ that event has probability zero, as the discussion after Lemma~\ref{lem:realized-cell} shows, and an indefinitely extensible family of questions is a different thing from a record that answers each of them rarely.
\noindent\textit{The status of the assumption.} {\rm (A1)} is an ordinary scientific-modeling claim about the cosmic weighting, contestable by a domain specialist attribute by attribute and the sort of thing evidence bears on. Causal separation of the events is relevant evidence for it but does not establish it: shared background parameters, or uncertainty about them, can make causally separated events statistically dependent, so an application must argue the bound for its own attributes. No assumption about the observer-conditioned weighting enters the test. An earlier calibration of a compounded multiplier did need such assumptions; the realized-cell bound supersedes it, and the long development keeps it as the record of how the bound was first calibrated \cite{Nexus}.

\emph{What survives without (A1).} A weaker bound is available using no independence assumption at all, only the marginal lower-tail super-uniformity of each $Q_i$, the observer-conditioned probability of the side of $\mathcal{A}_i$ the record realizes, written $A_i(X)$. By the union bound, with $\delta/m < \tfrac12$ so that the lower-tail bound of Theorem~\ref{thm:rarity} applies, $\Pr[\exists i:\, Q_i \le \delta/m] \le \delta$, so with probability at least $1-\delta$ every $Q_i$ exceeds $\delta/m$, and since $Q_i \le P(A_i(X))/P(H)$,
\begin{equation}
\label{eq:union-bound}
P(H) \;\le\; \frac{m}{\delta}\,\min_i P\big(A_i(X)\big) .
\end{equation}
On the reporting branch on which every realized side is the designated rare side this is $(m/\delta)\min_i p_i$; on other records the procedure reports the realized-side bound as displayed, in which each common-side realization contributes its own prior, the minimum still coming from a rarer realized side where there is one, and a small bound arising only from rare realized sides and within the error allowance. This is far weaker than the conjunction rule---it uses the rarest realized side rather than the product---but it is assumption-free past the single-attribute theorem, and it is what a reader who declines (A1) is left holding. A reader who keeps (A1) keeps the conjunction rule, $P(H) \le \prod_i p_i/\delta$ on the record that realizes every rare side, and with it the Rare Earth corollary, whose conjunction rule rests on {\rm (A1)} alone; the realized-cell bound, $P(H) \le (2^m/\delta)\,P(C(X))$ with coverage at least $1-\delta$ (Lemma~\ref{lem:realized-cell}), needs no {\rm (A1)} and reports on every record. 

The antecedent's clauses, then, in order of burden. The first is {\rm (A1)}, cosmic submultiplicativity over the predeclared family, for which causal separation is evidence rather than proof. The second is the record's length. The corollary asks for $m$ predeclared contingencies realized on their rare sides, with $m$ fixed by the target through $p^m \le \delta\varepsilon$; it does not ask that a biosphere's history supply infinitely many usable attributes, nor that the record answer every question on its rare side. How far a given record reaches is an empirical matter, and every application reads its level off a finite bound: $m$ attributes in hand, all realized on their rare sides, give $P(H) \le \prod_i p_i/\delta$ at confidence $1-\delta$, and a record realizing some common sides gives $P(H) \le (2^m/\delta)\, P(C(X))$ over its realized cell, a definite number either way. The third is the prior ceiling, and under the conjunction rule it is no longer a threshold: a binary partition's rare side has prior at most one half, so $p^m \le \delta\varepsilon$ is reached by $m$ alone. The threshold $p < \tfrac12$ survives only in the realized-cell rule, which also rejects on records realizing some common sides and pays $2^m$ for that reach. What a reader should press on is therefore the first clause, {\rm (A1)}, together with the predeclaration: those are the places the conclusion could fail without any of the mathematics being wrong. Nothing in the antecedent asks that any one event be rare to any particular degree; the rarity is manufactured by conjunction from ordinary contingency. The theorem specifies when a predeclared record rejects an abundance hypothesis; whether Earth's record satisfies those conditions remains an empirical and selection-modeling question.

\section{The theorem}
\label{sec:theorem}

\begin{theorem}[A statistical trilemma]
\label{thm:triad}
Under the hypotheses of Corollary~\ref{cor:rareearth} and the population alignment {\rm (Pop)} with density bound $C$, the three principles---\emph{(P1)} self-mediocrity; \emph{(P2)} cosmic-mediocrity; \emph{(P3$_\eta$)} the astrobiological Copernican principle with abundance threshold $\eta$---are tested together by a long enough predeclared rare-side record: (P1) and (P2) make the record a draw on which Corollary~\ref{cor:rareearth}'s test has level $\delta$, and (P3$_\eta$) with (Pop) entails the abundance hypothesis $P(H) \ge \eta/C$ that the test rejects. For every $\delta \in (0,1)$, a record realizing the rare sides of $m$ predeclared attributes with $p^m \le \delta\eta/C$ rejects $\mathrm{P3}_\eta$ at level $\delta$, and with it the triad's conjunction. The rejection is statistical, its level falling without bound in $m$, and not a logical contradiction: all three principles can be true while the test rejects, and the probability of that is the test's type-I error, at most $\delta$.
\end{theorem}

\begin{proof}
Three links. \emph{(i)} P1 and P2 jointly yield the Accordance Principle: that is the derivation of Section~\ref{sec:principle}, with the relocations of Section~\ref{sec:relocations} fixing the objects over which the mediocrity commitments range. \emph{(ii)} The principle bounds observer-abundance by the prior of any realized attribute, Theorem~\ref{thm:rarity}; conjunction drives the prior to zero, and the realized-cell bound gives the probability of the draw that lands there, Lemma~\ref{lem:realized-cell}; that half of the link is structural and is carried on this paper's page. The other half is the corollary's hypotheses, and Corollary~\ref{cor:rareearth} isolates what they are: cosmic submultiplicativity (A1) over a predeclared family, a prior ceiling on that family, which binary rare sides meet automatically under the conjunction rule, the strict $p < \tfrac12$ being needed only by the realized-cell rule---and a record realizing the rare sides of $m$ of its members, with $m$ fixed by the level. All three are assessable, at least in principle, from the record and the cosmic weighting; no assumption about the observer-conditioned weighting is needed. That is what stands between the principle and the Rare Earth hypothesis \cite{ward2000rare}, and the empirically assessable part of it is argued in \autoref{sec:rarity-bound} on this paper's page. \emph{(iii)} Rare Earth rejects P3, and the step deserves its own statement because the two claims are about populations that must be made to coincide. Theorem~\ref{thm:rarity} bounds $P(H)$, the cosmic-weighted probability that a natural history reaches observer-production; P3 asserts that observer-bearing environments like Earth's are common among planetary environments. These are the same proposition only under a condition, which is stated here rather than assumed:
\begin{description}[leftmargin=3.2em, style=nextline]
\item[(Pop) Population alignment.] The measure $\nu$ under which P3 is stated is absolutely continuous with respect to the cosmic weighting $\mu$ under which $P(H)$ is evaluated, with bounded density $d\nu/d\mu \le C$; the exact-alignment case, the two measures ranging over one population with one weighting, is $C = 1$.
\end{description}
Under (Pop), $\nu(H) \le C\,\mu(H)$, so if P3 is given the quantitative content it must have to be contradicted---$\mathrm{P3}_\eta$: $\nu(H) \ge \eta$ for a substantive abundance threshold $\eta > 0$, ``common'' meaning at least that---then the test rejects $\mathrm{P3}_\eta$ exactly when $B_m < \eta/C$, where $B_m$ is the upper confidence bound the test supplies: $\prod_i p_i/\delta$ under the conjunction rule of Corollary~\ref{cor:rareearth}, and $(2^m/\delta)\,P(C(X))$ under the realized-cell rule; at $C = 1$, ``$P(H)$ is small'' and ``observer-bearing environments are rare in the population P3 speaks of'' are one statement in two vocabularies. Without (Pop) the two may be reconciled by taking them over different populations, and a defender of P3 is entitled to that move---at the cost of saying which population their principle is about and why the cosmic weighting is not it. (Pop) is an additional modeling assumption, not merely a reading of P3. Calling both measures natural does not supply a density bound between the cosmic weighting over histories and the population measure over planetary environments, and the constant $C$ enters the conclusion quantitatively, the rejection threshold scaling as $\eta/C$. Every formulation of the astrobiological principle known to the author quantifies over planetary environments under some natural measure, which makes (Pop) the natural bridge but does not by itself justify its constant. The constant is supplied by the reference class, and the reference class by the structure of the cosmic weighting. Every planet is reachable from ours by perturbation: the Aethus describing Earth, edited attribute by attribute, yields that of Proxima Centauri~b, of a fictional world, or of a sterilized Earth \cite{Nexus}. Relocating from Earth to Proxima Centauri~b by perturbing the Aethus and relocating by moving through physical space reach the same description, so the physical cosmos is one rendition of the weighted space of alternatives, the \emph{Open Cosmos}, and the Open Cosmos is the only exhaustive reference class \cite{Nexus, Aethus}. Read over it, P3 is a claim about the cosmic weighting itself, and $C = 1$. What reachability establishes is the support of the population, that the Open Cosmos contains every planetary description; it does not establish that the population the surveys count carries the Open Cosmos's weights, two measures on the same states being free to weight them differently, so that the density bound $d\nu/d\mu \le C$ is a quantitatively justified input and reachability contributes nothing to its value; and a partial information state does not by itself specify the sampling measure or the granularity at which it is weighted, so the measure problem is relocated by the change of objects and not eliminated by it. $C = 1$ is therefore a stipulated alignment, the transparent case, and the class construction that follows is what computes $C$ from an explicit population. Read over any explicit class $S$ of planetary environments---rocky planets in habitable zones, say---weighted as the cosmic weighting weights its members, P3 says $\mu(H \mid S) \ge \eta$. Since $\mu(H) \ge \mu(H \cap S) = \mu(H \mid S)\,\mu(S)$, the population measure $\nu = \mu(\cdot \mid S)$ has density at most $1/\mu(S)$, so $C = 1/\mu(S)$: computed from the class, not assumed. It costs the test only logarithmically, $\log(1/\mu(S))/\log(1/p)$ further predeclared attributes, about three at $p = 0.2$ for a class carrying one percent of the cosmic weight. The one exit is a population weighted by something other than the cosmic weighting, and a defender who takes it must name that weighting and say why it, rather than the Open Cosmos, is what P3 is about. Two limits belong with the argument. Its motivating step, that any threshold of resemblance to Earth is artificial, is a sorites consideration; the explicit-class computation does not need it. And reading P3 over the planets that actually exist treats them as trials of planet-level processes, although they share one universe's laws and history and are viewed from a planet selected by its observers: the cosmic weighting is defined without correlation to one's own existence, which is what that reading must borrow, and the correlations of a single history bound how far actual frequencies can be trusted to estimate it. The exact-alignment case $C = 1$ is the one this paper illustrates. Chaining: if P1, P2, and P3 all held, links (i) and (ii) would deliver Rare Earth while P3 denies it. Hence the triad's conjunction is rejected at the level the record reaches, and every level short of certainty is reached by a long enough such record. The proof establishes the test only; which member to surrender is argued, not proved, in the next section.
\end{proof}

The conditional form of the theorem is deliberate, and it is worth saying what is conditional on what. Links (i) and (iii) are carried entirely on this paper's page: (i) is a two-line consequence of Bayes' theorem given the relocations, the structural half of (ii) is Theorem~\ref{thm:rarity}, and (iii) is a statistical rejection under the joint assumptions. No link imports; the empirical antecedent of (ii) is the weakest of the three and is argued where it is stated. A reader who rejects the antecedent---who holds that a natural history does not supply enough predeclarable rare-side contingencies for the level in question---keeps a conditional of independent interest, the mediocrity commitments and the Copernican principle standing or falling with that antecedent; a reader who grants it inherits the full trilemma.

A corollary formalizes the title. Call \emph{Copernican methodology} the commitment to typicality reasoning under a non-gerrymandered weighting---which is P1 and P2---and call the \emph{Copernican conclusion} the astrobiological thesis P3.

\begin{corollary}[The method against the principle]
\label{cor:method}
Conditional on link (ii) of Theorem~\ref{thm:triad}, and on the sampling model, prior specification, aggregation assumptions, and attribute-selection protocol of Section~\ref{sec:theorem}, a position that retains Copernican methodology faces a statistical rejection of the Copernican conclusion for the astrobiological case, at the confidence the theorem states; which premise to revise is a separate philosophical question, taken up in Section~\ref{sec:exits}.
\end{corollary}

\begin{proof}
Immediate: P1 and P2 jointly with link (ii) deliver Rare Earth at that confidence by links (i) and (ii), and Rare Earth at that confidence is a rejection of P3, not a refutation of it.
\end{proof}

\section{Pricing the exits}
\label{sec:exits}

% R845: a rejection implicates the conjunction with its auxiliary inputs; self-mediocrity granted by most parties (it
% read "At least one member fails" and "every party ... already grants").
A record meeting the criterion rejects a conjunction, the three theses with the sampling model, the prior ceilings, the attributes' selection and aggregation, and the population alignment, any of which a rejection can implicate. Granting the auxiliary inputs, at least one member of the triad fails at the confidence the record reaches, and the question is which surrender costs least. The test is proved; which member to surrender is not, and this section argues rather than demonstrates. What it offers is a comparison of costs, structural rather than a matter of empirical taste, and a reader who weighs the costs differently rejects the recommendation without rejecting the theorem.

\emph{Surrendering P1} proves too much. Self-mediocrity, in some operative form, is the move most parties to the observation-selection debates grant---it is what makes any anthropic inference an inference at all---and its outright denial licenses dismissing arbitrary observations with ``we are simply atypical,'' a solvent that dissolves the practitioner's own arguments before it reaches this paper's. Nor does keeping P1 commit one to vacuous self-typicality claims about particular objects. The constraint's typicality is substantive because it is taken across histories: the question is never whether \emph{this specific Earth} is typical of some class containing it---a comparison that is either trivial or ill-posed---but whether the observer-situation, specified attribute by attribute, is typical among observer-situations across the space of ways things could have gone. The specific-object version of the worry is dissolved, not answered, by the first relocation. One internal check belongs here as well: the Rare Earth conclusion itself does not offend P1, because P1's typicality is taken among observer-situations --- conditioned on $H$ --- and among observer-situations, occupancy of a live environment is universal however rare such environments are in the unconditioned weighting. The rarity lives unconditioned; the typicality claim lives conditioned; the relocation is what keeps the two from colliding.

\emph{Surrendering P2} costs the standard toolkit the reach of its own workhorse, and the cost can be exhibited rather than asserted. Survivorship-bias reasoning---the textbook correction every practitioner accepts---is precisely the present constraint collapsed to its coarsest limiting case. Restrict the state space to a single history's population, and let the observer-condition degenerate from the graded production-of-observers $H$ to a binary survival event $S$. The identity behind inequality~\eqref{eq:accordance} then reads
\begin{equation}
\label{eq:survivorship}
\frac{P(A \mid S)}{P(A' \mid S)} \;=\; \frac{P(S \mid A)}{P(S \mid A')} \cdot \frac{P(A)}{P(A')},
\end{equation}
the statement that observed frequency ratios among survivors equal prior ratios times survival-relevance ratios. Read as a correction---divide out the survival factor to recover the undistorted prior---this is the standard workhorse in its entirety; the Accordance Principle is the same identity read in the other direction, with $S$ regraded to $H$ and the single history reopened to the space of histories. Surrendering P2 leaves that correction intact, since it weighs one actual population and needs no weighting over histories. What it forfeits is the correction's extension from one history's population to the space of histories, which is the step the Rare Earth inference takes. The cosmic weighting is therefore not an exotic addition but the familiar correction's weighting carried from one history to all of them; a tailored weighting, by contrast, reintroduces the reference-class disease as a free parameter, since every tailoring is a choice of class by other means.

\emph{Surrendering P3} is, on this comparison, the least costly, and---this is the paper's central structural point---it is the exit recommended by the very method P3's defenders trust. The Copernican principle's support was never a body of exoplanetary evidence; it was a method: the Copernican demotion of privileged position, generalized by Gott into an explicit typicality inference \cite{gott1993implications}. The present framework is that method generalized further---from positions in one history to attributes across histories---and the generalization, fed the natural-historical record, is what produces the Rare Earth conclusion \cite{Nexus, ward2000rare}. An independent modern route deserves one sentence of convergence: propagating realistic uncertainties through the standard population estimates already yields substantial probability that we are alone \cite{sandberg2018dissolving} --- a different method, sharing no premise with the present one, arriving in the same region. The convergence is noted, not imported; link (ii) does not lean on it. The method, run to completion, turns against the principle it was thought to support.

\textit{What remains of P3, and the Great Filter.} What remains of P3 once its methodological support defects is habit; and the surrender has a further consequence worth stating in the filter vocabulary \cite{hanson1998greatfilter}. The Great Filter framing asks where, along the path from dead matter to visible civilization, the improbability that explains the silent sky is concentrated---with the uncomfortable possibility that it lies ahead of us. Surrendering P3 does not settle the behind-or-ahead question, and the framework is structurally silent on it: the Accordance constraint conditions on observer-production, which our own case has already secured, so no bound it issues bears on what becomes of observers afterward. What the surrender does is narrower and still consequential---past rarity of environments and histories like ours becomes sufficient to explain the silent sky, which removes one abductive reason to posit a filter ahead; it does not predict that no such filter exists, and a reader who takes ``the filter's mass lies behind us'' from this paper has taken more than it gives. The companion's program then does more than a relocation invites: rather than asking \emph{which} past attribute carries the mass, it shows by a substitution property that the constraint can be written at any pair of stages, an earlier and a later, so that every adjacent pair carries its own filter result and the question ``where is the filter'' presupposes a location the constraint denies \cite{Nexus}. The mass is distributed across the rungs, not concentrated at one; that quantification is the companion's, not this paper's.

\textit{The recommended resolution.} The recommended resolution is therefore to retain the mediocrity commitments and surrender the astrobiological Copernican principle---holding, with full awareness of the irony, that Copernican reasoning itself is what demotes the Copernican conclusion.

\section{Objections and replies}
\label{sec:objections}

One discipline governs every reply below. The test rejects a conjunction, of abundance, of the sampling model, of the transport assumptions and of the prior bounds, and rejecting a conjunction does not say which conjunct failed; that a realized rare feature must have received compensating observer-relevance is one horn of that trilemma, the one the paper's reading takes, and not what the rejection by itself establishes.

\emph{First objection: the cosmic weighting merely relocates the reference-class problem.} The worry: calling one weighting ``cosmic'' is itself a choice, so the disease diagnosed in Section~\ref{sec:relocations} has been renamed, not cured. The reply distinguishes two problems the tradition runs together. Choosing a \emph{class} of observers within a completed history is the anthropic-specific disease, and it genuinely cannot close, because any boundary drawn through a completed history's population is as defensible as any other. Choosing a \emph{prior measure} over ways things can go is the ordinary burden of all Bayesian reasoning, shared with every scientific inference and disciplined by the usual means---symmetry, naturalness, robustness of conclusions under variation. On this paper's reading, the two relocations convert the first problem into the second. The constraint's conclusions, moreover, are ratios, and inherit robustness from that form: they survive any reweighting that does not differentially target the compared pair.

\emph{Second objection: this is a fine-tuning argument in disguise.} It is not, and the difference is the object of reasoning. Fine-tuning arguments reason over fundamental constants, with observer-selection applied to an ensemble of complete cosmologies---the completed-timeline unit again, one level up \cite{bostrom2002anthropic}. The present constraint binds attribute by attribute \emph{within} natural history, over partial states, and issues in local, quantitative claims about the observer-relevance of particular realized attributes. The framework is consistent with any resolution of the fine-tuning debate and lends support to none of its parties; the only shared element is the word ``typicality,'' relocated here to a different object.

\emph{Third objection: the selection-effect truism already covers this.} The truism---observers necessarily find themselves where observers can exist, so improbable-looking situations call for no further explanation---runs in exactly one direction: it \emph{discharges} realized improbability, given observers. The Accordance Principle runs in the other: it \emph{converts} realized improbability into forced, graded claims of observer-relevance, with magnitudes that propagate into further reasoning. The truism is recovered as the degenerate limit of the constraint in which the forced magnitudes are discarded; retaining them is the entire content of the method, and is what the truism, by construction, cannot supply.

\emph{Fourth objection: multiplicity, and the choice of alternative.} Two worries share a root. First, examine enough realized attributes and some will be the a priori unlikely member of their pair by sheer number --- everyone is atypical somewhere --- so the principle threatens to force spurious relevance attributions wherever one looks hard enough. Second, the alternative $A'$ looks like a free parameter: pair the realized attribute against a suitably chosen rival and the bound says what the analyst wishes.

\textit{Reply to the first: the exceedance form is the correction.}\label{sec:multiplicity} The reply to the first is that the exceedance form~\eqref{eq:exceedance} is itself the multiplicity guard, and the guard is quantitative: each interrogated pair admits a $\kappa$-fold violation with probability at most $1/(1+\kappa)$, so among $n$ interrogated pairs the expected number of $\kappa$-fold violations is at most $n/(1+\kappa)$, by linearity and nothing else. Two quantities are distinct, since the weaker is easily mistaken for the stronger: the \emph{expected number} of exceedances and the \emph{probability of at least one}. The paper claims the stronger, and it comes free. Writing $N$ for the number of exceedances among the $n$ pairs, $N$ is non-negative and integer-valued, so Markov's inequality gives $\Pr[N \ge 1] \le \mathbb{E}[N] \le n/(1+\kappa)$, with no independence among the pairs assumed. The guard is therefore familywise---the probability that \emph{any} interrogated pair exhibits a $\kappa$-fold violation is at most $n/(1+\kappa)$, not merely the average count---and what it requires is that $n$ count the whole admissible family, which is the whole of the burden the selective-search objection imposes. Admissibility has a definition: an \emph{admissible family} $\mathcal{F}_0 = \{A_1, \ldots, A_n\}$ is declared \emph{before} inspecting which members are surprising, and the theorem applies simultaneously over $\mathcal{F}_0$; an analyst who chooses attributes after looking must count the full candidate family from which the reported attribute was selected, not the reported attribute alone. The canonical complement settles $A_i' = A_i^{c}$ once $A_i$ is chosen; it does not settle which $A_i$ belong to $\mathcal{F}_0$, and that admissibility question is the paper's, stated as a condition on the use of the theorem and not as a footnote. Forced inferences are those whose ratios stand clear of that noise floor --- a discipline of the same multiple-comparison kind every such inference owes, with the correction here supplied by the principle's own tail profile rather than imposed from outside. This discipline is not merely owed but observed: the companion conducts the multiple-comparison audit explicitly for its applications \cite{Nexus}, and the claims here are at the principle grade.

\textit{Reply to the second: the alternative is canonical.} The reply to the second is that $A'$ is not analyst-supplied but canonical: in the principle's boxed form the alternative is the complement --- $P(A') = 1 - P(A)$ --- so there is nothing to choose. The founding illustration, properly parsed, then runs on two canonical pairs rather than one: $A_1$, the dinosaurs arose (complement: never arose --- the canonical complement), and $A_2$, having arisen, they perished (complement: persisted); the coupling of these two pairs is the linked-set structure --- improbabilities in twos, never alone --- that the companion's founding argument turns on \cite{Nexus}, so the canonicality reply and the illustration reinforce rather than strain each other. Finer alternatives require refining the crossroads into an explicit partition, and refinement answers to the partition-modal form of the constraint, not to the analyst's convenience. What remains free is which attributes to interrogate, and that freedom, by the first reply, moves the noise floor and nothing else.

\textit{One concession.} One concession belongs in the open: $n$ must count the attributes implicitly \emph{scanned}, not merely those explicitly interrogated --- striking attributes get interrogated because a scan flagged them, a post-selection the noise-floor reply covers only when $n$ counts the scan's true breadth.

\section{What this paper does not claim}
\label{sec:fence}

The fence is a boundary, not a disparagement: the natural-historical program built on the principle, its conjectures and its quantitative bounds, is the companion's \cite{Nexus}, and this paper stops at the test, the requirements of its application, and the one consequence named next.

\subsection{Coupled improbabilities, and where this paper stops}
\label{sec:coupled}

One consequence of the principle is worth naming, because it is what a follow-up development takes as its object and because it marks this paper's boundary precisely.

% R845: the closure in the companion's form, the powerset of attribute-label pairs under fixed rules; an extensive
% operator stabilizes on a finite lattice without monotonicity, which gives leastness (it read, on partial
% assignments, "where F is not monotone the iteration need not converge", with monotonicity from positive coupling).
The principle does not terminate where it is applied. Attributing observer-relevance to one realized attribute constrains the others the constraint touches: if $A$ is credited with the relevance its improbability demands, then attributes on which $A$'s relevance depends inherit commitments of their own, and the record's improbabilities are therefore the \emph{closure} of its root commitments under a propagation step rather than a tally of independent factors. Two labels or none, in the simplest case: a prerequisite and what it is prerequisite to stand or fall together. The closure can be given as an operator, in the companion's form \cite{Nexus}. On the powerset of the finite set of attribute--label pairs, a complete lattice under inclusion, consistency handled by the companion's fiberwise clause and not by the domain, with the rules finitely many and fixed with their tails, priced in the principle's form, before iteration, let
\begin{equation}
\label{eq:closure-operator}
F(\mathsf{S}) \;=\; \mathsf{S} \;\cup\; \{\, \ell : \text{some rule with its premises in } \mathsf{S} \text{ concludes } \ell \,\}.
\end{equation}
Being extensive, $F$ stabilizes from any root set, monotone or not; being monotone, since a rule applicable at $\mathsf{S}$ stays applicable on every superset, it has a least fixed point above the roots by the Knaster--Tarski theorem, reached where that iteration stabilizes. Forcing read as conditioning would not be monotone, since conditioning on more can lower a probability; fixed rules keep the closure deductive, its failure priced once by the sum of the tails. This paper stops at that existence. What the closure \emph{is}---whether it concentrates on an isolated cell or spreads across a plateau---is the subject of a follow-up development, which adds hypotheses this paper does not make: a high-contrast hypothesis, that for moderately correlated attributes exchanging an attribute for its alternative moves observer-productivity by large and effectively random orders of magnitude, together with the structural assumptions needed to turn local contrast into global mass collapse---a product structure with cumulative suppression, and the licenses its own ledger inventories. Under those hypotheses the closure is shown to be the record's optimal cell. Nothing here depends on it; the rarity conclusion above is reached before it is assumed, and what it adds is not the rarity but its structure. The present paper is the foundation that development rests on, and this subsection is the seam.

\section{Conclusion}
\label{sec:conclusion}

% R845: the conclusion consolidated at the strength the test supports: the test rejects a conjunction and which member
% goes is argued (it read "whether that exit is taken is what the test decides"); the reference-class problem an artifact
% "on this reading"; the commitments' price argued, not forced; the repeated summary and the second and third paragraphs
% merged.
One program, one move, two wings: remove the classical bookkeeping, and the paradox is recast. The physics wing removes it from determinacy; the present wing removes it from typicality, and the two relocations---reference class across histories, partial states in place of completed timelines---are that removal made formal. What the removal buys, in this paper's scope, is a derived, local, quantitative typicality constraint; a proposed diagnosis of the reference-class problem as a completion artifact, conditional on the admissibility principle stated with it; a rarity bound and the calibrated test of observer abundance it yields; and a theorem: under the stated sampling and prior assumptions, a record satisfying the rejection criterion supplies a level-$\delta$ test against the abundance hypothesis that self-mediocrity, cosmic-mediocrity, and the astrobiological Copernican principle jointly entail. The argued resolution keeps the mediocrity commitments and would surrender the Copernican principle, on the ground that the Copernican method itself, generalized, is what tells against the principle; the test can reject the conjunction once its inputs are supplied, a rejection implicating the auxiliary inputs as well as the theses, and which member goes is argued, not decided by the test. Until the inputs are supplied, which the empirical-status table shows they are not yet, the Rare Earth conclusion remains a proposed application.

What follows is the paper's diagnosis of the anthropic literature, not a consequence of the derivation. The dominant tradition inherited a unit---the completed history, tagged for whether it houses observers---from the classical picture of a determinate world, and built its debates on that unit: which reference class to draw within the completed history, how to weight observers across completed worlds. Relocate the reasoning to partial states across possible histories, and the typicality principle becomes \emph{derived} rather than postulated, falling out of two mediocrity commitments by a single application of Bayes' theorem; the reference-class problem becomes, on this reading, an \emph{artifact} of that unit, replaced by an explicit measure-and-sampling problem in which reachability fixes the support of the measure and not its values; and the principle becomes \emph{quantitative and local}, issuing bounded attribute-wise claims where the tradition issued global verdicts. Whether the replacement settles the original dispute remains philosophical work. The two root commitments, the partial state as the object of reasoning and mediocrity as an inequality between two states, are not forced by logic; the paper argues that declining them has a price, the reference-class problem returning in its received form (Section~\ref{sec:relocations}) or the Copernican method losing its content, since ``the observer is not special'' is the only thing it ever said.

The natural-historical program is deliberately not this paper's. Everything further---the coupling of improbabilities that the founding argument first exhibited, the conjectures and quantitative bounds, the paleontological predictions---belongs to the companion and to a sibling extraction, where each claim is graded in its own register. One program, one move, two wings---and this is the anthropic wing's foundation stone, cut to stand on its own.


\paragraph{A note on provenance.} Two remarks, offered for orientation rather than as load-bearing argument. The first concerns the principle's origin. The Accordance Principle was, in the framework's own record, first stated in an unusually concrete setting: a question in natural history about whether the end-Cretaceous extinction was a prerequisite for the eventual evolution of observers like us \cite{Nexus}. The original form of the argument asked why, given that we find ourselves on the branch of history in which dinosaurs evolved and then perished, we should not instead have expected the apparently simpler branch in which they never arose at all---and inferred, from the demand that our own situation be typical rather than exceptional, that the never-dinosaurs branch must be unfavorable to observers for a reason the argument forces but does not itself supply. That early argument already exhibits, in miniature, everything this paper has stated in general form: the constraint binding a particular attribute pair, the direction running from realized improbability to forced observer-relevance, and the refusal to complete either branch into a total history before reasoning over it. It also exhibits the propagation the companion's program is built on---a typicality principle does not terminate where it is applied, but forces correlated attributions elsewhere in the record---though that development is the companion's to carry, with \autoref{sec:coupled} offering only its preliminary walkthrough.

The second remark is intellectual-historical. The apparatus was generalized from the Doomsday argument \cite{gott1993implications, leslie1996end} in isolation from the anthropic-principle literature, and the principle was disentangled from its founding application in 2021 so that it could stand as a general statement, which is where Nexic reasoning as a formalism begins; the dated documentary chain recording both steps lives in the companion \cite{Nexus}. The point is not priority but structure: an apparatus built by extending the observation-selection literature would inherit that literature's unit of reasoning, and this one demonstrably was not and does not.

\section*{Provenance}

% R823 (Fable 5.1, Oct 18 2026, author-directed): the origin of Nexic reasoning, dated from the long paper's record (its
% subsection on the biographical origin of the Accordance Principle; the September 18, 2021 and December 2021 entries;
% the 2021 provenance of the High-Contrast Earth Hypothesis, the drops and the Hateg hypothesis; the December 2024 deposit).
% R842 (Opus 5.5, Oct 22 2026), at the author's direction: the Hateg hypothesis dated by his journal entry of March 15,
% 2022 (the note had listed it among the 2021 originations, "the Hateg hypothesis and the Fluke problem of December 2021").
\noindent Nexic reasoning, whose foundation this paper states, was first developed in 2021. The Accordance Principle had been reached by mid-2020 inside the dinosaurian argument of \autoref{sec:coupled}, whose first written record is a video script of June 9, 2020; in 2020 the principle and that argument were one thing, and in 2021 the author deliberately separated them, stating the principle as a general logical statement and building the calculus on it, with the Optimal- and Perfect-Earth Conjectures (June--July 2021), the Theorem of Aimless Optimation (September 18, 2021), the High-Contrast Earth Hypothesis and the Imperfection Magnitude Drops, and the Fluke problem of December 2021 all dating from that year; the Ha\c{t}eg hypothesis followed, in March 2022. That formalism of 2021 and 2022 took the principle at face value; the Aethic reasoning of 2022 and 2023 motivated it from the Aethic ontology instead, and the long Nexic paper, first deposited in December 2024, records both \cite{Nexus}. The present paper was first posted in July 2026.

\noindent The two-stage reading of the Accordance Principle, with \emph{teleological discharge} as its second stage, follows the author's note of 2026-10-02 (2:19 PM ET), quoted verbatim in the long paper; the naming is the author's.

\printbibliography

\end{document}